\documentclass[11pt]{article}

\usepackage[margin=1in]{geometry}
\usepackage{amsmath, amssymb}
\usepackage{mathrsfs}
\usepackage{bm}
\usepackage{xcolor}
\usepackage{fancyhdr}
\usepackage{booktabs}
\usepackage{multirow}
\usepackage{array}
\usepackage{url}

\pagestyle{plain}


\begin{document}

% =========== Page 1: wei_page1.tex ===========

% ---------- Journal nameplate ----------
\noindent
{\footnotesize Contents lists available at ScienceDirect}\\[2pt]
{\Large\bfseries Applied Mathematical Modelling}\\[2pt]
{\footnotesize journal homepage: \texttt{www.elsevier.com/locate/apm}}

\vspace{0.8em}
\hrule
\vspace{1em}

% ---------- Title ----------
{\large\bfseries
Nondestructively identifying the mechanical behavior of soft
tissues using surface deformation with an explicit inverse
approach\par}

\vspace{0.8em}

% ---------- Authors ----------
\noindent
Yue Mei\textsuperscript{a,b},
Dongmei Zhao\textsuperscript{a,b},
Changjiang Xiao\textsuperscript{a,b},
Zhi Sun\textsuperscript{a,b},
Weisheng Zhang\textsuperscript{a,b,*},
Xu Guo\textsuperscript{a,b,*}

\vspace{0.5em}

% ---------- Affiliations ----------
\noindent
{\footnotesize
\textsuperscript{a}\,\emph{State Key Laboratory of Structural Analysis
for Industrial Equipment, Department of Engineering Mechanics, Dalian
University of Technology, Dalian 116023, PR China}\\[2pt]
\textsuperscript{b}\,\emph{International Research Center for
Computational Mechanics, Dalian University of Technology, Dalian
116023, PR China}\par}

\vspace{1em}
\hrule
\vspace{1em}

% ---------- Article info / Abstract block (two columns) ----------
\noindent
\begin{minipage}[t]{0.30\textwidth}
\textbf{A R T I C L E\ \ I N F O}\\[0.8em]
\textbf{Keywords:}\\
Inverse problem\\
Layered structures\\
Inclusion embedded structures\\
Surface deformation\\
Biological tissue
\end{minipage}
\hfill
\begin{minipage}[t]{0.66\textwidth}
\textbf{A B S T R A C T}\\[0.8em]
Identifying the spatial variation of stiffness properties in soft
tissues nondestructively, with limited surface measurements, poses
significant challenges. In this paper, we present a novel explicit
inverse approach designed to characterize the nonhomogeneous elastic
property distribution of soft tissues using only surface displacement
datasets. In contrast to the prevalent implicit inverse approach,
which focuses on optimizing the elastic properties of individual
pixels, our proposed method optimizes the geometric parameters of
deformable and movable components, as well as shear moduli of each
component. As a result, the proposed approach requires far fewer
optimization variables, streamlining the process. Numerical tests
conducted in this study demonstrate the superiority of the explicit
inverse method over the implicit inverse method, providing
much-improved reconstructed results. In particular for a ring
structure, while the average relative error of reconstruction using
the implicit inverse method can exceed 40\,\%, the explicit inverse
method achieves a remarkable average relative error of only 5\,\%.
Given that surface displacements are easily measurable, the
integration of our proposed explicit inverse method with low-cost
imaging techniques shows great potential in accurately mapping the
elastic property distribution of biological tissues.
\end{minipage}

\vspace{1.2em}
\hrule
\vspace{0.8em}

% ---------- Introduction ----------
\section*{Introduction}

\noindent
Biological tissues often display significant heterogeneity. This
diversity is evident in various tissues such as skin [1--3], blood
vessels [4,5], and corneas [6,7], among others, which exhibit
multi-layered structures. These layers offer a variety of
functionalities, thus contributing to the complexity of biological
systems. Interestingly, the emergence of cancerous tumors introduces
an additional level of heterogeneity to these tissues. These tumors
typically exhibit greater stiffness compared to their surrounding
healthy tissues due to pathological alterations [8]. This stark
contrast in mechanical properties can serve as a potential indicator
of malignancy. It's also worth noting that the local mechanical
properties of soft tissues are not static; they fluctuate in response
to various factors such as age

% ---------- Bottom-of-page footer (corresponding-author block) ----------
\vfill
\hrule
\vspace{0.4em}
\noindent
{\footnotesize
* Corresponding authors at: State Key Laboratory of Structural
Analysis for Industrial Equipment, Department of Engineering
Mechanics, Dalian University of Technology, Dalian 116023, PR China.\\[2pt]
\textit{E-mail addresses:}
\texttt{weishengzhang@dlut.edu.cn} (W.~Zhang),
\texttt{guoxu@dlut.edu.cn} (X.~Guo).\\[4pt]
\url{https://doi.org/10.1016/j.apm.2024.05.028}\\[2pt]
Received 10 October 2023; Received in revised form 21 May 2024;
Accepted 22 May 2024\\[2pt]
Available online 31 May 2024\\[2pt]
0307-904X/\textcopyright{} 2024 Elsevier Inc. All rights are reserved,
including those for text and data mining, AI training, and similar
technologies.\par}

\newpage

% =========== Page 2: wei_page2.tex ===========

\noindent
\ldots{}[9--11], dietary habits [12], and the presence of pathological
conditions [13,14]. Therefore, understanding the nonhomogeneous
biomechanical behavior of these tissues is of paramount importance,
particularly in a non-destructive manner. Such knowledge can inform
medical diagnostics and treatments, potentially leading to better
patient outcomes.

Considerable effort has been devoted into research aimed at
identifying the nonhomogeneous biomechanical properties of biological
tissues. Studies [15--18] have demonstrated the feasibility of
estimating the mechanical properties of each layer in multi-layered
biological tissues, provided that the thickness of each layer is known
beforehand. However, in practice, determining the thickness of each
layer prior to addressing the identification problem can be a
challenging task. Furthermore, the interface between neighboring
layers is seldom flat, a characteristic particularly pronounced in
biological tissues. As such, a comprehensive solution requires
knowledge of both the topology of the interface between the layers
and the biomechanical properties of each layer. Additionally, the
detection of tumors embedded within normal tissues introduces another
layer of complexity. To accurately diagnose and plan treatment for
these conditions, it is crucial to determine not only the elastic
modulus values of these tumors but also their topology [19].

A general strategy for tackling this challenge involves identifying
the spatial variation of material properties, a process that
necessitates solving a large-scale optimization problem to optimize
the elastic properties of every pixel in the domain of interest,
which can be referred to as the implicit inverse method [20--23]. The
implicit inverse method has been successfully employed in tumor
detection using real clinical data [8,24]. However, this implicit
inverse method has its limitations. It not only demands significant
computational resources but is also prone to severe uniqueness
issues, which can hamper the reliability of the results. Given that
biological tissues are often known a priori to exhibit either layered
structures or inclusion-embedded structures, it's feasible to
explicitly and simultaneously identify the topology of each
nonhomogeneous region and its associated material properties. This
approach bypasses some of the computational hurdles of the implicit
method. Another inherent drawback of the implicit inverse method lies
in the quality of the reconstruction results, which can be
compromised if the full-field displacement field cannot be measured
with high accuracy. Studies [25,26] that employed the implicit
inverse method to reconstruct the nonhomogeneous elastic property
distribution of solids using surface deformation have demonstrated
this issue. While these methods are capable of visualizing the
nonhomogeneous region, the quality of the reconstruction is highly
sensitive to noise. This sensitivity arises from a disparity between
the large number of unknown optimization variables in the implicit
inverse method and the sparsity of the deformation information on the
surface. In contrast, the explicit inverse method significantly
reduces the total number of optimization variables [27--29],
potentially improving the quality of the reconstructed results. This
strategy holds promise for enhancing the accuracy of nonhomogeneous
biomechanical property identification in biological tissues.

This paper will present an explicit inverse method that utilizes
surface deformation measurements to identify the mechanical behavior
of two common nonhomogeneous structures: layered structures and
inclusion embedded structures. The organization of this work is as
follows: In the \textbf{Methods} Section, we provide a detailed
explanation of the mathematical background underlying the explicit
inverse approach. The \textbf{Results} Section presents a comparative
study between the explicit and implicit inverse approaches, using
various numerical examples. We discuss the obtained results in the
\textbf{Discussion} Section and conclude the paper with a summary in
the \textbf{Conclusion} Section.

\section*{Methods}

\subsection*{\textit{Forward problem}}

\noindent
Given the low elastic modulus of soft tissues, this study will
consider finite deformation in the analysis. The strong form of the
nonlinear elastic problem is stated as follows:
\begin{equation}
\left\{
\begin{aligned}
\mathrm{Div}(\mathbf{P}) &= \mathbf{0} && \text{in } \Omega \\
\mathbf{u}               &= \mathbf{g} && \text{on } \Gamma_g \\
\mathbf{P}\cdot\mathbf{N} &= \mathbf{T} && \text{on } \Gamma_T
\end{aligned}
\right.
\end{equation}
where $\mathbf{P} = \mathbf{F}\mathbf{S}$ is the first
Piola--Kirchhoff stress tensor where $\mathbf{F}$ is the deformation
gradient and $\mathbf{S}$ is the second Piola--Kirchhoff stress
tensor. $\Omega$ is the problem domain at the reference
configuration. $\mathbf{g}$ is the prescribed displacement on the
displacement boundary $\Gamma_g$. $\mathbf{T}$ is the prescribed
traction on the traction boundary $\Gamma_T$. The following
restrictions apply to $\Gamma_g$ and $\Gamma_T$: $\Gamma_g \cup
\Gamma_T = \partial\Omega$ and $\Gamma_g \cap \Gamma_T = \phi$. We
assume the soft tissue of interest obeys the incompressible and
neo-Hookean law:
\begin{equation}
\mathbf{S} = \mu\!\left(\mathbf{I} - \mathbf{C}^{-1}\right)
             + p\,\mathbf{C}^{-1}
\end{equation}
where $\mathbf{C} = \mathbf{F}^{\mathsf{T}}\mathbf{F}$ is the
Cauchy--Green tensor. $p$ is the hydrostatic pressure.

The weak form of the nonlinear elastic forward problem is stated as:
Find $[\mathbf{u}^h, p^h] \in \mathscr{M}^h \times \mathscr{P}^h$
such that
\begin{equation}
\mathscr{A}\!\left(\mathbf{w}^h,\mathbf{u}^h,q^h,p^h\right)
\;=\;
\int_{\Gamma_T}\! \mathbf{w}^h\!\cdot\mathbf{T}\,\mathrm{d}\Gamma,
\qquad
\forall\,[\mathbf{w}^h, q^h] \in \ell^h \times \mathscr{P}^h
\end{equation}
\begin{equation}
\mathscr{A}\!\left(\mathbf{w}^h,\mathbf{u}^h,q^h,p^h\right)
\;\equiv\;
\int_\Omega \nabla\mathbf{w}:\mathbf{P}\,\mathrm{d}\Omega
\;+\;
\int_\Omega q\,(J-1)\,\mathrm{d}\Omega
\;-\;
\sum_{i=1}^{NE}\!
 \Bigl(\tau\,\mathrm{div}(\mathbf{F}^h\mathbf{S}^h),\;
       (\mathbf{F}^h)^{-\mathsf{T}}\nabla q^h\Bigr)_{\Omega_i}
\end{equation}

\newpage

% =========== Page 3: wei_page3.tex ===========
\setcounter{equation}{4}

\noindent
where the superscript $h$ denotes the finite element subspace.
$\mathscr{M}^h$, $\ell^h$ and $\mathscr{P}^h$ are the finite-element
subspaces of the function space
$\mathscr{M} \equiv \{\mathbf{u} \mid u_i \in H^1(\Omega);\,
u_i = g_i \text{ on } \Gamma_g\}$,
$\ell \equiv \{\mathbf{w} \mid w_i \in H^1(\Omega);\,
w_i = 0 \text{ on } \Gamma_g\}$ and
$\mathscr{P} \subseteq L^2(\Omega)$, respectively. The third term is
the stabilized term to address the volumetric locking induced by the
incompressibility where $NE$ is the total number of elements in the
discretized domain. $\tau = \dfrac{\alpha h^{2}}{2\mu}$ represents the
stabilization factor, where $h$ is the characteristic element length.
$\alpha$ is set to $1.0$ [30].

\subsection*{\textit{Topological description of nonhomogeneous elastic property distribution}}

\noindent
To describe the topology of the nonhomogeneous elastic property
distribution, we introduce the B-spline function to express a curve:
\begin{equation}
\mathbf{C}(t) = \sum_{i=0}^{n+1} N_{i,p}(t)\,\mathbf{P}_i,
\qquad (a \leq t \leq b)
\end{equation}
where $N_{i,p}(t)$ is the $p$-th order B-spline basis function
generated by a knot vector $\{t_0, t_1, \ldots, t_{n+2+p}\}$.
$\mathbf{P}_i\,(i = 0, \ldots, n+1)$ is the $i$-th control point.

The function $N_{i,p}(t)$ can be determined by Eq.~(6):
\begin{equation}
\begin{cases}
N_{i,0}(t) =
\begin{cases}
1, & \text{if } t_i \leq t \leq t_{i+1} \\
0, & \text{otherwise}
\end{cases}
& (i = 0, \ldots, n+1;\ p = 0) \\[3.5ex]
N_{i,p}(t) = \dfrac{t - t_i}{t_{i+p} - t_i}\,N_{i,p-1}(t)
+ \dfrac{t_{i+p+1} - t}{t_{i+p+1} - t_{i+1}}\,N_{i+1,p-1}(t),
& (i = 0, \ldots, n+1;\ p \geq 1)
\end{cases}
\end{equation}

To capture various types of nonhomogeneous elastic property
distributions, we employ B-splines to approximate both open and
closed curves, as depicted in Fig.~1. The open curve is utilized to
represent the interface shape between adjacent layers in a layered
composite, while the closed curve represents the shape of inclusions
embedded within a soft solid. To avoid the self-intersection of the
closed B-spline curve, we define the control points of the closed
B-spline curve in Eq.~(7):
\begin{equation}
\mathbf{P}_i =
\begin{cases}
\left(
x_o + d_i \sin\!\left(i\theta - \dfrac{\theta}{2}\right),\;
y_o + d_i \cos\!\left(i\theta - \dfrac{\theta}{2}\right)
\right),
& \left(\theta = \dfrac{2\pi}{n};\ i = 1, \ldots, n\right) \\[3ex]
\mathbf{P}_0 = \mathbf{P}_{n+1} = \dfrac{\mathbf{P}_1 + \mathbf{P}_n}{2}
\end{cases}
\end{equation}
where $(x_o, y_o)$ is the location of the center of the closed
B-spline curve, $d_i$ is the distance between the $i$-th control
point and the center of the closed curve.

With the curve expressed by the B-spline function, we can define
each nonhomogeneous region in the problem domain by

% ---------- Figure 1 placeholder ----------
\vspace{1em}
\begin{center}
\fbox{\parbox{0.92\textwidth}{\centering\vspace{1em}
\textbf{[Figure 1 — placeholder]}\\[0.8em]
\textbf{(a) Open curve}: eight control points
$\mathbf{P}_0, \mathbf{P}_1, \ldots, \mathbf{P}_7$ connected by dashed
blue segments, with a smooth orange B-spline curve passing through
the polygon.\\[0.4em]
\textbf{(b) Closed curve}: nine control points
$\mathbf{P}_0 = \mathbf{P}_9, \mathbf{P}_1, \ldots, \mathbf{P}_8$
distributed around a center $(x_o, y_o)$, with radial distances
$d_1, d_2, \ldots, d_8$ and angular spacing $\theta = 2\pi/n$.
The closed B-spline curve (orange) is inscribed in the dashed
control polygon, with $\mathbf{P}_0 = \mathbf{P}_9 = (\mathbf{P}_1 + \mathbf{P}_8)/2$.
\vspace{1em}}}
\end{center}
\vspace{0.4em}
\begin{center}
\textbf{Fig. 1.} Two examples of the curve expressed by B-spline functions.
\end{center}

\newpage

% =========== Page 4: wei_page4.tex ===========
\setcounter{equation}{7}

\noindent
introducing the following topological description function (TDF):
\begin{equation}
\begin{cases}
\phi^i(\mathbf{x}) < 0 & \text{if } \mathbf{x} \in \Omega^i \\
\phi^i(\mathbf{x}) = 0 & \text{if } \mathbf{x} \in \partial\Omega^i \\
\phi^i(\mathbf{x}) > 0 & \text{if } \mathbf{x} \in \Omega \setminus \Omega^i
\end{cases}
\end{equation}

For the layered structure, $\phi^i$ is the TDF of the $i$-th
interface and $\Omega^i$ is the region occupied from the curve $C^i$
and $C^0$. The mathematical expression of TDF for a layered structure
can be defined as:
\begin{equation}
\phi^i(\mathbf{x}) = r(t) - r^i(t)
\end{equation}
where $r^i(t)$ is the distance between the curve $C^i$ and $C^0$ at
any arbitrary point on the domain. To this end, the $m$-th layer of
the structure ($m > 1$) can be defined by $\phi^{m-1}(\mathbf{x}) > 0$
and $\phi^{m}(\mathbf{x}) \leq 0$, except for the first layer
$\phi^{1}(\mathbf{x}) < 0$ and last layer $\phi^{n}(\mathbf{x}) > 0$.
To this end, A bilayer structure can be expressed by the proposed
method as shown in Fig.~2. Note that we assume that each layer is
homogeneous and did not consider spatial variations within a layer.

For the domain with inclusion embedded, $\phi^i$ is the TDF of the
$i$-th component and $\Omega^i$ is the region occupied by the $i$-th
component. For an inclusion embedded composite, the mathematical
expression of TDF is:
\begin{equation}
\phi^i(\mathbf{x}) = \sqrt{\bigl(x - x_o^i\bigr)^{2}
                           + \bigl(y - y_o^i\bigr)^{2}}
                     - r^i\!\left(\theta^i\right)
\end{equation}
where $\theta^i$ is the angle from the point $(x, y)$ to the center
of the $i$-th component $\left(x_o^i, y_o^i\right)$ in polar
coordinate system. $r^i(\theta^i)$ is the distance from the point
with the angle of $\theta^i$ on $\partial\Omega^i$ to the center
$\left(x_o^i, y_o^i\right)$ as shown in Fig.~3.

With the defined TDF, the spatial variation of the shear modulus of
a layered composite can be defined based on the Boolean operators
[31]:
\begin{equation}
\mu(\mathbf{x}) = \mu^{0}\bigl(1 - H(\phi^{1}(\mathbf{x}))\bigr)
+ \sum_{i=1}^{n-1}\!\bigl[\mu^{i}\, H(\phi^{i}(\mathbf{x}))
                           \bigl(1 - H(\phi^{i+1}(\mathbf{x}))\bigr)\bigr]
+ \mu^{n}\, H(\phi^{n}(\mathbf{x}))
\end{equation}
where $\mu^{0}$ is the shear modulus of background. $\mu^{n}$ is the
shear modulus of the last layer. $H(z)$ is the Heaviside function.
If there is no interface curve $C^i$, $\mu(\mathbf{x}) = \mu^{0}$ is
a homogeneous domain. However, the TDF definitions of composite with
embedded inclusions are different, that is:
\begin{equation}
\boldsymbol{\mu}\!\left(\boldsymbol{x}\right) =
\boldsymbol{\mu}^{0} \prod_{i=1}^{n}
  \mathbf{H}\!\left(\boldsymbol{\phi}^{i}(\boldsymbol{x})\right)
+ \sum_{i=1}^{n}\!\bigl(\boldsymbol{\mu}^{i}
  \bigl(1 - \mathbf{H}(\boldsymbol{\phi}^{i}(\boldsymbol{x}))\bigr)\bigr)
\end{equation}

As the Heaviside function is a discontinuous function, we will
utilize the following function to approximate the Heaviside function
so that we can evaluate the spatial gradient of the Heaviside
function:
\begin{equation}
H(z) = 0.5 \times \left(1 + \tanh\!\left(\dfrac{z}{\varepsilon}\right)\right)
\end{equation}
where $\varepsilon$ is a small constant. With the decrease of
$\varepsilon$, the approximation of the Heaviside function is closer
to the corresponding analytical expression (see Fig.~4). However, a
smaller $\varepsilon$ might lead to an extremely large gradient at
$z$ approaching to zero. Thus, the small constant $\varepsilon$
should be carefully chosen.

\subsection*{\textit{Inverse problem}}

\noindent
After discussing the TDF, the inverse problem can be stated as:
Given $n_{\text{meas}}$ surface measured displacement fields, find
the optimal geometric vectors $\mathbf{r}^1, \mathbf{r}^2, \ldots,
\mathbf{r}^m$ and the shear moduli $\mu^0, \mu^1, \ldots, \mu^m$,
such that the following objective function is minimized:
\begin{equation}
\Pi = \dfrac{1}{2} \sum_{i_{\text{meas}}=1}^{n_{\text{meas}}}
\int \left|\mathbf{u}_i^{\text{com}} - \mathbf{u}_i^{\text{meas}}\right|^{2}
\, d\partial\Omega_{i_{\text{meas}}} + \text{Reg}(\mu)
\end{equation}

% ---------- Figure 2 placeholder ----------
\vspace{1em}
\begin{center}
\fbox{\parbox{0.92\textwidth}{\centering\vspace{1em}
\textbf{[Figure 2 — placeholder]}\\[0.8em]
Bilayer structure illustrating the TDF construction. The bottom
layer $\mu^{0}$ (teal) is bounded below by the reference curve
$C^{0}$; the top layer $\mu^{1}$ (pink) is bounded above by the
interface curve $C^{1}$. Six control points
$\mathbf{P}^{1}_{1}, \mathbf{P}^{1}_{2}, \ldots, \mathbf{P}^{1}_{6}$
define the B-spline interface $C^{1}$ (orange). Radial distances
$r^{1}_{1}, r^{1}_{2}, \ldots, r^{1}_{6}$ connect each control point
back to $C^{0}$, parameterised along $t$.
\vspace{1em}}}
\end{center}
\vspace{0.4em}
\begin{center}
\textbf{Fig. 2.} An example of the nonhomogeneous shear modulus
distribution of a bilayer structure.
\end{center}

\newpage

% =========== Page 5: wei_page5.tex ===========
\setcounter{equation}{14}

% ---------- Figure 3 placeholder ----------
\noindent
\begin{center}
\fbox{\parbox{0.92\textwidth}{\centering\vspace{1em}
\textbf{[Figure 3 — placeholder]}\\[0.8em]
Square light-teal domain with background shear modulus $\mu^{0}$,
containing an irregular pink inclusion region with shear modulus
$\mu^{1}$ and boundary $\partial\Omega^{1}$. Eight control points
$\mathbf{P}^{1}_{1}, \mathbf{P}^{1}_{2}, \ldots, \mathbf{P}^{1}_{8}$
are distributed around the boundary, connected to each other by
dashed blue segments (the control polygon) and smoothly interpolated
by an orange closed B-spline curve. Eight radial distances
$d^{1}_{1}, d^{1}_{2}, \ldots, d^{1}_{8}$ connect the inclusion
center $(x^{1}_{0}, y^{1}_{0})$ to each control point, illustrating
the polar parameterization used in Eq.~(7) and Eq.~(10).
\vspace{1em}}}
\end{center}
\vspace{0.3em}
\begin{center}
\textbf{Fig. 3.} The nonhomogeneous shear modulus distribution of an
inclusion embedded composite.
\end{center}

\vspace{0.8em}

% ---------- Figure 4 placeholder ----------
\noindent
\begin{center}
\fbox{\parbox{0.92\textwidth}{\centering\vspace{1em}
\textbf{[Figure 4 — placeholder]}\\[0.8em]
Graph of $H(z) = 0.5 \times (1 + \tanh(z/\varepsilon))$ (the smoothed
Heaviside from Eq.~(13)) plotted over $z \in [-0.5, 0.5]$, with
$H(z) \in [0, 1]$ on the vertical axis. Four curves show the effect
of decreasing $\varepsilon$:\\[0.3em]
\begin{tabular}{ll}
\textcolor{red}{$\varepsilon = 0.01$} & sharpest transition (nearly a true step) \\
\textcolor{green!60!black}{$\varepsilon = 0.05$} & moderate transition \\
\textcolor{blue}{$\varepsilon = 0.10$} & smoother transition \\
\textcolor{orange}{$\varepsilon = 0.15$} & smoothest transition
\end{tabular}\\[0.3em]
All four curves pass through $(0, 0.5)$. As $\varepsilon \to 0$ the
approximation converges to the true Heaviside; as $\varepsilon$ grows
the transition band widens.
\vspace{1em}}}
\end{center}
\vspace{0.3em}
\begin{center}
\textbf{Fig. 4.} An approximation of Heaviside function.
\end{center}

\vspace{1em}

% ---------- Eq. 15 ----------
\noindent
\begin{equation}
\text{Reg}(\mu) = \dfrac{1}{2}\,\alpha \int_{\Omega}
\sqrt{|\nabla\mu|^{2} + c^{2}}\, d\Omega
\end{equation}

\vspace{0.3em}

% ---------- Prose after Eq. 15 ----------
\noindent
For the layered structure,
$\mathbf{r}^i = \left[r^i_1(t_1),\, r^i_2(t_2),\, \cdots,\, r^i_n(t_n)\right]$
is the geometric parameter vector representing the distance between
the $i$-th interface and $C^0$ at the location $t_k$. For the
structure with inclusions embedded, the geometric parameter vector
$\mathbf{r}^i = \left[r^i_1,\, r^i_2,\, \cdots,\, r^i_n\right]
= \left[x^i_0,\, y^i_0,\, d^i_1,\, \cdots,\, d^i_{n-2}\right]$
is consisting of the coordinates of the center and distances between
the control points to the center of $i$-th component. In Eq.~(14),
the discrepancy between the computed displacement field and measured
displacement field on partial surface $\partial\Omega_{i_{\text{meas}}}$
is minimized in $L$-$2$ norm. The computed displacement field is
obtained by solving a forward problem using the finite element method
[32]. Due to the influence of noise in measurements, we introduce a
total variation diminishing regularization (TVD) term (the second
term in Eq.~(14)), and the specific form is given by Eq.~(15).
$\alpha$ is the regularization factor which controls the significance
of the regularization on the objective function. $c$ is a small
constant to avoid singularities when evaluating the gradient of the
regularization term with respect to optimization variables. It should
be noted that there exist alternative forms for the regularization
term apart from the domain integral forms. Nevertheless, we have
chosen to adhere to the use of domain integrals due to practical
considerations. In particular, the explicit inverse solver has been
developed based on the established framework of the existing
implicit inverse solver. This continuity in methodology ensures
consistency and facilitates seamless integration between the two
solvers [33]. The specific value of the regularization factor is
chosen based on the smoothness criteria [22].

\newpage

% =========== Page 6: wei_page6.tex ===========
\setcounter{equation}{15}

\noindent
The optimization problem is solved by the limited
Broyden--Fletcher--Goldfarb--Shanno (L-BFGS) method which requires
the objective function and the gradients of the objective function
with respect to the optimization variables. The gradients of the
objective function with respect to the optimization variables are
obtained by the adjoint method which has been thoroughly discussed
in [27]. The flow chart of the procedure to solve the inverse
problem is shown in Fig.~5.

To test the feasibility of the proposed inverse scheme, we adopt the
displacement datasets generated by a forward finite element
simulation with a target shear modulus distribution. The white
Gaussian random is also added into the simulated surface displacement
field. The noise level is defined by
$\sqrt{\displaystyle\sum_{k=1}^{N_{\text{node}}}
\left(u_k - u_k^{\text{noise}}\right)^{\!2}\bigg/
\sum_{k=1}^{N_{\text{node}}} (u_k)^{2}} \times 100\,\%$,
where $u_k$ and $u_k^{\text{noise}}$ are the simulated and noisy
displacement of $k$-th node, $N_{\text{node}}$ is the total number of
nodes in interest domain. In comparison, we also solve inverse
problems using the prevalent implicit inverse method based on the
adjoint method where the shear moduli are nodally defined [22]. To
quantitatively analyze the performance of the proposed method, we
define the average relative error by
$\dfrac{1}{N_{\text{node}}}
\displaystyle\sum_{k=1}^{N_{\text{node}}}
\left|\left(\mu_k - \mu_k^{\text{target}}\right)
\big/ \mu_k^{\text{target}}\right| \times 100\,\%$\,,
where $\mu_k$ is the recovered shear modulus value of $k$-th node
and $\mu_k^{\text{target}}$ is the target shear modulus value of
$k$-th node.

\subsection*{\textit{Gradient computation}}

\noindent
To efficiently compute the gradient of the objective function with
respect to the optimization variable, we have adopted the adjoint
method, which involves solving the associated adjoint equation. The
adjoint equation is given by [34,35]:
\begin{equation}
\mathscr{B}\!\left(\mathbf{w}^h, q^h, \delta\mathbf{u}^h, \delta p^h
;\, \mathbf{u}^h, p^h\right)
= -\bigl(\delta\mathbf{u}^h,\,
         \mathbf{u}^h - \mathbf{u}^{m}\bigr)_{\partial\Omega},
\qquad
\forall\,[\delta\mathbf{u}^h, \delta p^h]
\in \mathscr{M}^h \times \mathscr{P}^h
\end{equation}
where
$\mathscr{B}(\mathbf{w}^h, q^h, \delta\mathbf{u}^h, \delta p^h;\,
\mathbf{u}^h, p^h) = \frac{d}{d\varepsilon}
\mathscr{A}^{*}\!(\mathbf{w}^h, q^h, \mathbf{u}^h
+ \varepsilon\,\delta\mathbf{u}^h, p^h
+ \varepsilon\,\delta p^h)\big|_{\varepsilon=0}$
is the linearization of $\mathscr{A}^{*}$. With the acquired state
variables $\mathbf{u}^h, p^h$, we can solve for the dual variables
$\mathbf{w}^h, q^h$ using Eq.~(16). To this end, the gradient of the
objective function with respect to the shear modulus:
\begin{equation}
\delta\Pi = D_{\mu}\Pi\,\delta\mu = \mathscr{C}
+ \alpha \int_{\Omega}
\dfrac{\nabla\mu \cdot \nabla\delta\mu}
      {\sqrt{|\nabla\mu|^{2} + c^{2}}}\, d\Omega
\end{equation}
$\mathscr{C} = \frac{d}{d\varepsilon}
\mathscr{A}^{*}\!(\mathbf{w}^h, q^h, \mathbf{u}^h, p^h,
\mu + \varepsilon\,\delta\mu)\big|_{\varepsilon=0}$.
This can be used to efficiently compute the gradient of the objective
function with respect to

% ---------- Figure 5 placeholder ----------
\vspace{1em}
\begin{center}
\fbox{\parbox{0.92\textwidth}{\centering\vspace{1em}
\textbf{[Figure 5 — placeholder: optimization flowchart]}\\[0.8em]
\begin{tabular}{l}
\fbox{Give initial guess: $\mathbf{r}^1, \mathbf{r}^2, \ldots,
      \mathbf{r}^m$ and $\mu^0, \mu^1, \ldots, \mu^m$} \\[0.6em]
\hspace{2em}$\downarrow$\quad \textit{Finite Element Method} \\[0.4em]
\fbox{Solve forward problem $\to$ computed displacement
      $\mathbf{u}^{\text{com}}$} \\[0.6em]
\hspace{2em}$\downarrow$\quad \textit{Adjoint Method} \\[0.4em]
\fbox{Evaluate $\Pi$ and its gradients w.r.t.\ optimization
      variables} \\[0.6em]
\hspace{2em}$\downarrow$ \\[0.4em]
\fbox{Is $\Pi$ minimized?}
$\xrightarrow{\textit{No}}$
\fbox{Update variables (\textit{L-BFGS})}
$\to$ loop back $\uparrow$ \\[0.6em]
\hspace{2em}$\downarrow$\quad \textit{Yes} \\[0.4em]
\fbox{Stop and output reconstruction results}
\end{tabular}
\vspace{1em}}}
\end{center}
\vspace{0.3em}
\begin{center}
\textbf{Fig. 5.} The flow chart of the proposed explicit inverse method.
\end{center}

\newpage

% =========== Page 7: wei_page7.tex ===========
\setcounter{equation}{17}

\noindent
nodal or elemental shear moduli. If the shear modulus is nodally
defined, we can use the same interpolation function with the
displacement field. In this case, the shear modulus can be expressed
as:
\begin{equation}
\mu \approx \mu^{h} = \sum_{i=1}^{N} N_i\, \mu_i
\end{equation}
where $N$ is the total number of nodes in the finite element model
and $N_i$ is the finite element shape function. $\mu_i$ is the nodal
shear modulus.

In the explicit inverse method, since shear modulus distribution is a
function of geometric parameters and the shear modulus in each
region, the gradient of the objective function with respect to the
optimization variables can be calculated using the chain rule. To
this end, we should seek $\delta\mu$ corresponding to a change in any
optimization variable from $\beta$ to $\beta + \delta\beta$:
\begin{equation}
\delta\mu = \frac{\partial\mu}{\partial\beta}\,\delta\beta
\end{equation}
Thereby, comparing to the implicit inverse method, we need to
evaluate the $\partial\mu / \partial\beta$.

\medskip
\noindent
\textbf{(a) Layered structure}

\medskip
\noindent
The evaluation of the gradient of geometric optimization variables in
a layered structure with $n + 1$ layers should be discussed
considering three cases. For the $m$-th geometric optimization
variable $r^{1}_{m}$ at the first interface $C^{1}$,
$\frac{\partial\mu}{\partial r^{1}_{m}}$ is given by:
\begin{equation}
\frac{\partial\mu}{\partial r^{1}_{m}}
= \left[-\mu^{0} + \mu^{1}\!\left(1 - H\!\left(\phi^{2}(\mathbf{x})\right)\right)\right]
  \frac{\partial H\!\left(\phi^{1}(\mathbf{x})\right)}{\partial\phi^{1}(\mathbf{x})}\,
  \frac{\partial\phi^{1}(\mathbf{x})}{\partial r^{1}_{m}}
\end{equation}

For the $m$-th geometric optimization variable $r^{n}_{m}$ at the
last interface $C^{n}$,
$\frac{\partial\mu}{\partial r^{n}_{m}}$ is given by:
\begin{equation}
\frac{\partial\mu}{\partial r^{n}_{m}}
= \left(-\mu^{n-1}\, H\!\left(\phi^{n-1}(\mathbf{x})\right) + \mu^{n}\right)
  \frac{\partial H\!\left(\phi^{n}(\mathbf{x})\right)}{\partial\phi^{n}(\mathbf{x})}\,
  \frac{\partial\phi^{n}(\mathbf{x})}{\partial r^{n}_{m}}
\end{equation}

For the $m$-th geometric optimization variable at the intermediate
interface $C^{i}\,(2 \leq i \leq (n-1))$,
$\frac{\partial\mu}{\partial r^{i}_{m}}$ is given by:
\begin{equation}
\frac{\partial\mu}{\partial r^{i}_{m}}
= \left[-\mu^{i-1}\, H\!\left(\phi^{i-1}(\mathbf{x})\right)
        + \mu^{i}\!\left(1 - H\!\left(\phi^{i+1}(\mathbf{x})\right)\right)\right]
  \frac{\partial H\!\left(\phi^{i}(\mathbf{x})\right)}{\partial\phi^{i}(\mathbf{x})}\,
  \frac{\partial\phi^{i}(\mathbf{x})}{\partial r^{i}_{m}}
\end{equation}

In Eqs.~(20--22),
$\frac{\partial H(\phi^{j}(\mathbf{x}))}{\partial\phi^{j}(\mathbf{x})}$
$(1 \leq j \leq n)$ can be evaluated by Eq.~(13) as the Heaviside
function is approximated by a continuous and differentiable function.
$\frac{\partial\phi^{j}(\mathbf{x})}{\partial r^{j}_{m}}$ can be
evaluated by the finite difference method.

For the shear modulus of the background $\mu^{0}$,
$\frac{\partial\mu}{\partial\mu^{0}}$ is given by:
\begin{equation}
\frac{\partial\mu}{\partial\mu^{0}} = 1 - H\!\left(\phi^{1}(\mathbf{x})\right)
\end{equation}

For the shear modulus of the last layer $\mu^{n}$, we have
\begin{equation}
\frac{\partial\mu}{\partial\mu^{n}} = H(\phi^{n}(\mathbf{x}))
\end{equation}

For the shear modulus of the layer $\mu^{i}\,(1 \leq i \leq (n-1))$,
we have
\begin{equation}
\frac{\partial\mu}{\partial\mu^{i}}
= H\!\left(\phi^{i}(\mathbf{x})\right)
  \left(1 - H\!\left(\phi^{i+1}(\mathbf{x})\right)\right)
\end{equation}

\medskip
\noindent
\textbf{(b) Domain with inclusion embedded}

\medskip
\noindent
For each $m$-th geometric optimization variables of $j$-th component
$r^{j}_{m}$,
$\frac{\partial\mu}{\partial r^{j}_{m}}$ is given by:
\begin{equation}
\frac{\partial\mu}{\partial r^{j}_{m}}
= \left[\mu^{0}\!\prod_{\substack{i=1 \\ i \neq j}}^{n}
         H\!\left(\phi^{i}(\mathbf{x})\right)
  - \mu^{j}\right]
  \frac{\partial H\!\left(\phi^{j}(\mathbf{x})\right)}{\partial\phi^{j}(\mathbf{x})}\,
  \frac{\partial\phi^{j}(\mathbf{x})}{\partial r^{j}_{m}}
\end{equation}

For the shear modulus of the background $\mu^{0}$, we have:

% [page ends here — continues on page 8]

\newpage

% =========== Page 8: wei_page8.tex ===========
\setcounter{equation}{26}

\begin{equation}
\frac{\partial\mu}{\partial\mu^{0}}
= \prod_{i=1}^{n} H\!\left(\phi^{i}(\mathbf{x})\right)
\end{equation}

For the shear modulus $\mu^{j}$ of the $j$-th component,
$\frac{\partial\mu}{\partial\mu^{j}}$ is given by:
\begin{equation}
\frac{\partial\mu}{\partial\mu^{j}}
= 1 - H\!\left(\phi^{j}(\mathbf{x})\right)
\end{equation}

\section*{Results}

\subsection*{\textit{Case 1: A layered rectangular structure}}

\noindent
In the first numerical example, we analyze a layered rectangular
structure, illustrated in Fig.~6. The dimensions of this structure
are 1.0\,mm $\times$ 0.2\,mm. To discretize the problem domain, we
utilize 7200 linear triangular elements. In the forward finite
element simulation, a concentrated force is applied to the top edge,
while the bottom edge was restricted the vertical motion and the
center point is fully fixed. In the inverse problem, our focus lies
in utilizing displacement datasets collected from the top edge as
inputs for the inverse scheme. Acknowledging that relying on a single
surface displacement dataset may not produce optimal reconstruction
results, we change the location of the concentrated force and
incorporate multiple surface displacement datasets to solve the
inverse problem. The target shear modulus distributions for one (the
homogenous case), two, and three layers are depicted in Fig.~6(a),
(b), and (c), respectively. The shear modulus values of each layer in
the structure are of a similar order to that of skin tissue [36].
Considering the uncertainty regarding the total number of layers, we
initially assume a configuration with six layers for the inverse
problem, as depicted in Fig.~7(a). For the explicit inverse scheme,
we utilize 10 control points to interpolate each interface. Thus, the
total number of optimization variables is 56 including 50 geometric
variables and 6 shear moduli.

From Fig.~7, it can be seen that both the shape of the interfaces
between the neighboring layers and the shear modulus values of each
layer varies and finally reach to the target shear modulus
distribution (Sample 1). Note that it takes only 52 iterations for
convergence. Figs.~8--10 are the reconstructed shear modulus
distributions for Samples 1--3 using explicit and implicit inverse
approaches for the noise free cases, respectively. In the implicit
inverse method, the total number of optimization variables is
linearly proportional to the mesh size. In this case, the total
number of optimization variables is 3691, which is equal to the total
number of nodes. The reconstructed results demonstrate that both
inverse methods can yield decent reconstruction with noise free
cases. Further, the implicit method fails to recover the
heterogeneity at the upper corner compared with the results by the
explicit inverse method. Since the total number of optimization
variables is remarkably reduced in the explicit inverse method, the
total number of minimization iterations is reduced. The total number
of minimization iterations for the explicit inverse method is
approximately 10 times less than that required for the implicit
inverse method. Furthermore, the shear modulus reconstruction
obtained through the implicit

% ---------- Figure 6 placeholder ----------
\vspace{1em}
\begin{center}
\fbox{\parbox{0.92\textwidth}{\centering\vspace{1em}
\textbf{[Figure 6 — placeholder]}\\[0.8em]
Three layered rectangular structures (1.0\,mm $\times$ 0.2\,mm),
each showing the target shear modulus (SM) distribution in MPa
with a jet-style colorbar on the right.\\[0.4em]
\textbf{(a) Sample 1}: homogeneous, single layer, SM $= 3.00$\,MPa
(uniform blue). Downward arrows on top edge indicate concentrated
force locations. Red text ``displacement measurement'' with
leftward arrows on left side.\\[0.3em]
\textbf{(b) Sample 2}: two layers. Bottom layer SM $\approx 3$\,MPa
(blue), top layer SM $\approx 10$\,MPa (red). Green circles on top
boundary, red circles on bottom boundary mark displacement
measurement points.\\[0.3em]
\textbf{(c) Sample 3}: three layers (blue--red--blue sandwich). SM
ranges 1.00--10.00\,MPa. Same boundary marker pattern as (b).\\[0.3em]
All three samples share the same mesh (7200 triangular elements) and
boundary conditions (top-edge force, bottom-edge vertical
constraint, center point fixed).
\vspace{1em}}}
\end{center}
\vspace{0.3em}
\begin{center}
\textbf{Fig. 6.} The target shear modulus distribution of three
samples with different number of layers (Unit: MPa). SM stands for
shear modulus.
\end{center}

\newpage

% =========== Page 9: wei_page9.tex ===========

% ---------- Figure 7 placeholder ----------
\noindent
\begin{center}
\fbox{\parbox{0.92\textwidth}{\centering\vspace{1em}
\textbf{[Figure 7 — placeholder]}\\[0.8em]
Six-panel grid showing the optimization trajectory for the
explicit inverse scheme on \textbf{Sample 1} (the homogeneous case,
target $\text{SM} = 3.00$\,MPa), starting from a 6-layer initial
guess. Each panel is a rectangular shear-modulus distribution
($1.0$\,mm $\times 0.2$\,mm) with a jet colorbar.\\[0.4em]
\begin{tabular}{@{}ll@{}}
\textbf{(a) Initial guess} & uniform blue, SM $= 1.00$\,MPa everywhere \\
\textbf{(b) Iteration 5} & SM range 1.27--1.95, layered artifacts emerging \\
\textbf{(c) Iteration 12} & SM range 1.50--3.11, layered structure developed \\
\textbf{(d) Iteration 24} & SM range 1.63--3.03, approaching target \\
\textbf{(e) Iteration 48} & SM range 2.97--3.00, nearly converged \\
\textbf{(f) Iteration 52 (Final)} & uniform red, SM $= 3.00$\,MPa everywhere
\end{tabular}\\[0.4em]
The trajectory shows the 6-layer parameterisation collapsing to a
uniform solution over 52 iterations --- because the target is
homogeneous, the optimiser correctly drives all six layers to the
same shear modulus.
\vspace{1em}}}
\end{center}
\vspace{0.3em}
\begin{center}
\textbf{Fig. 7.} The optimization process with six layers initially
utilizing the explicit inverse scheme for Sample 1 (Without noise,
Unit: MPa).
\end{center}

\vspace{0.8em}

% ---------- Figure 8 placeholder ----------
\noindent
\begin{center}
\fbox{\parbox{0.92\textwidth}{\centering\vspace{1em}
\textbf{[Figure 8 — placeholder]}\\[0.8em]
Six-panel side-by-side comparison of the
\textbf{Explicit Inverse Method} (left column) vs.\
\textbf{Implicit Inverse Method} (right column) on Sample 1
(homogeneous, target SM $= 3.00$\,MPa), with 3, 6, and 9 surface
displacement datasets as inputs.\\[0.4em]
\textbf{Explicit (left column)}: all three rows show a perfectly
uniform shear-modulus field at SM $= 3.00$\,MPa (uniform blue,
colorbar pinned 3.00--3.00). The explicit parameterisation
recovers the homogeneous target exactly.\\[0.4em]
\textbf{Implicit (right column)}: all three rows show tiny spatial
artifacts around the true value of 3.00\,MPa --- (d) 3 fields: SM
range 2.9995--3.0001; (e) 6 fields: SM range 2.9992--3.0001;
(f) 9 fields: SM range 3.0000--3.0002. The color blobs are
visible even though the numerical error is on the order of
$10^{-4}$--$10^{-3}$\,MPa.\\[0.4em]
\textbf{Takeaway}: the explicit method produces a perfectly uniform
field on the homogeneous case; the implicit method is quantitatively
accurate but shows residual spatial noise patterns that decrease
(slightly) as more displacement datasets are incorporated.
\vspace{1em}}}
\end{center}
\vspace{0.3em}
\begin{center}
\textbf{Fig. 8.} The reconstructed results with varying 3, 6 and 9
surface displacement datasets corresponding to Sample 1 (Without
noise, Unit: MPa).
\end{center}

\vspace{1em}

% ---------- Continuation prose (from page 8) + noise-test paragraph ----------
\noindent
inverse method shows a slight spatial variation, which can be
enhanced by incorporating more surface datasets.

After adding 1\,\% noise to the surface displacement datasets, we
observed that the shear modulus distributions were accurately mapped
for both inverse methods in the homogeneous case (refer to Fig.~11).
The average relative error for both implicit and explicit inverse
methods was less than 0.2\,\% (see Table~1). However, when it came
to Sample 2 and 3 (see Fig.~12, Fig.~13), the interface between
neighboring layers in the implicit inverse method's results was not
as clear as the target. As the total number of layers increases, the
challenges in solving the inverse problems escalate accordingly.
Consequently, the relative error for Sample 2 and Sample 3 exhibits
a lower quality compared to Sample 1 (see Table~1). Besides, the
average relative error for the implicit inverse method was

% [page ends here — continues on page 10]

\newpage

% =========== Page 10: wei_page10.tex ===========

% ---------- Figure 9 placeholder ----------
\noindent
\begin{center}
\fbox{\parbox{0.92\textwidth}{\centering\vspace{1em}
\textbf{[Figure 9 — placeholder]}\\[0.8em]
Six-panel side-by-side comparison of the \textbf{Explicit Inverse
Method} (left column) vs.\ \textbf{Implicit Inverse Method} (right
column) on \textbf{Sample 2} (the two-layer target from Fig.~6b,
nominal SM values $\mu^0 \approx 3$\,MPa, $\mu^1 \approx 10$\,MPa),
using 3, 6, and 9 surface displacement datasets as inputs.\\[0.4em]
\textbf{Explicit (left column)}: bilayer structure recovered with a
clean interface. Colorbar ranges:
(a) 3 fields, SM $\in [1.00, 11.49]$\,MPa;
(b) 6 fields, SM $\in [1.00, 11.63]$\,MPa;
(c) 9 fields, SM $\in [1.00, 11.47]$\,MPa.\\[0.3em]
\textbf{Implicit (right column)}: bilayer structure recovered but
with a visibly rougher interface between the two layers; more
artifacts at the top edge. Colorbar ranges:
(d) 3 fields, SM $\in [1.00, 10.95]$\,MPa;
(e) 6 fields, SM $\in [1.00, 10.77]$\,MPa;
(f) 9 fields, SM $\in [1.00, 11.13]$\,MPa.\\[0.3em]
\textbf{Takeaway}: both methods recover the two-layer structure;
the explicit method produces a cleaner interface, while the implicit
method's interface has more small-scale roughness. Slight overshoot
in the top-layer SM values (colorbar max above the nominal 10\,MPa)
is visible in all six panels.
\vspace{1em}}}
\end{center}
\vspace{0.3em}
\begin{center}
\textbf{Fig. 9.} The reconstructed results with varying 3, 6 and 9
surface displacement datasets corresponding to Sample 2 (Without
noise, Unit: MPa).
\end{center}

\vspace{0.8em}

% ---------- Figure 10 placeholder ----------
\noindent
\begin{center}
\fbox{\parbox{0.92\textwidth}{\centering\vspace{1em}
\textbf{[Figure 10 — placeholder]}\\[0.8em]
Six-panel side-by-side comparison of the \textbf{Explicit Inverse
Method} (left column) vs.\ \textbf{Implicit Inverse Method} (right
column) on \textbf{Sample 3} (the three-layer target from Fig.~6c:
blue bottom layer $\mu^0 \approx 3$\,MPa, green middle layer
$\mu^1 \approx 6$\,MPa, red top layer $\mu^2 \approx 10$\,MPa),
using 3, 6, and 9 surface displacement datasets as inputs.\\[0.4em]
\textbf{Explicit (left column)}: three-layer stacked structure
recovered, with the middle (green) layer visibly distinct from the
top (red) and bottom (blue) layers. Colorbar ranges:
(a) 3 fields, SM $\in [1.00, 10.49]$\,MPa;
(b) 6 fields, SM $\in [1.00, 10.35]$\,MPa;
(c) 9 fields, SM $\in [1.00, 10.44]$\,MPa.\\[0.3em]
\textbf{Implicit (right column)}: three-layer structure recovered
but with fuzzier interfaces between adjacent layers. Colorbar
ranges:
(d) 3 fields, SM $\in [1.00, 10.11]$\,MPa;
(e) 6 fields, SM $\in [1.00, 10.26]$\,MPa;
(f) 9 fields, SM $\in [1.00, 10.19]$\,MPa.\\[0.3em]
\textbf{Takeaway}: as layer count increases (1 $\to$ 2 $\to$ 3),
the quality gap between explicit and implicit widens. Sample 3's
three-layer structure is harder for the implicit method to resolve
cleanly than Sample 2's two-layer structure. More displacement
datasets (3 $\to$ 6 $\to$ 9) improve both methods but do not close
the gap.
\vspace{1em}}}
\end{center}
\vspace{0.3em}
\begin{center}
\textbf{Fig. 10.} The reconstructed results with 3, 6 and 9 surface
displacement datasets corresponding to Sample 3 (Without noise,
Unit: MPa).
\end{center}

\vspace{1em}

% ---------- Continuation prose (from page 9) + 5% noise paragraph ----------
\noindent
slightly larger than that for the explicit inverse method (see
Table~1). Furthermore, it's worth noting that the quality of the
reconstruction achieved by both inverse methods improved with an
increase in the total number of measurements. By utilizing 9
boundary displacement measurements, the average relative error
could be reduced to less than 10\,\%. However, it should be
mentioned that each layer was not accurately mapped as the
homogeneous case when using the implicit inverse method.

Additionally, we increased the noise level to 5\,\% and observed
that the reconstructed results obtained by the explicit inverse
method still outperform those obtained by the implicit inverse
method (see Fig.~14, Fig.~15, Fig.~16). For the explicit inverse
method, the resulting relative error in the case with 5\,\% noise
is generally larger than that in the case with 1\,\% noise (see
Table~1). However, even with a 5\,\% noise level, the relative
error for the explicit inverse method remains lower than that for
the implicit inverse method

% [page ends here — continues on page 11]

\newpage

% =========== Page 11: wei_page11.tex ===========

% ---------- Figure 11 placeholder ----------
\noindent
\begin{center}
\fbox{\parbox{0.92\textwidth}{\centering\vspace{1em}
\textbf{[Figure 11 --- placeholder]}\\[0.8em]
Six-panel side-by-side comparison of the \textbf{Explicit Inverse
Method} (left column) vs.\ \textbf{Implicit Inverse Method} (right
column) on \textbf{Sample 1} (homogeneous, target SM $= 3.00$\,MPa),
with \textbf{1\,\% noise} added to the measured surface
displacements. Three rows correspond to 3, 6, and 9 displacement
datasets.\\[0.4em]
\textbf{Explicit (left column)}:
(a) 3 fields: SM $\in [3.0013, 3.0013]$\,MPa, uniform red;
(b) 6 fields: SM $\in [3.0000, 3.0001]$\,MPa, visible green/yellow
patches emerge (non-monotone: noisier than (a) or (c));
(c) 9 fields: SM $\in [2.9997, 2.9997]$\,MPa, uniform red.\\[0.3em]
\textbf{Implicit (right column)}:
(d) 3 fields: SM $\in [2.9981, 3.0083]$\,MPa, visible blue patch
along the left edge;
(e) 6 fields: SM $\in [2.9978, 3.0006]$\,MPa, small central blue
blob;
(f) 9 fields: SM $\in [2.9995, 2.9998]$\,MPa, small blue patches
remain near the edges.\\[0.3em]
\textbf{Takeaway}: under 1\,\% noise, both methods stay
quantitatively close to the true 3.00\,MPa value ($\le 0.01$\,MPa
deviation). The explicit method's 6-field result is the noisiest
explicit panel on this page --- an artifact not discussed in the
surrounding prose.
\vspace{1em}}}
\end{center}
\vspace{0.3em}
\begin{center}
\textbf{Fig. 11.} The reconstructed results with varying 3, 6 and 9
surface displacement datasets corresponding to Sample 1 (1\,\% noise
is introduced into the measured data, Unit: MPa).
\end{center}

\vspace{1em}

% ---------- Table 1 ----------
\noindent\textbf{Table 1}\\[0.2em]
{\small Comparison of the average relative error in the shear
modulus distribution obtained by implicit inverse and explicit
inverse methods for the layered structure.}
\vspace{0.3em}

\begin{center}
\small
\begin{tabular}{@{}ll ccc ccc@{}}
\toprule
\multicolumn{2}{l}{Inverse method}
  & \multicolumn{3}{c}{Explicit inverse method}
  & \multicolumn{3}{c}{Implicit inverse method} \\
\cmidrule(lr){3-5}\cmidrule(lr){6-8}
\multicolumn{2}{l}{Number of surface displacement datasets}
  & 3 & 6 & 9 & 3 & 6 & 9 \\
\midrule
\multirow{3}{*}{1\,\% noise}
  & Sample 1 & 0.04\,\%  & 0.00\,\%  & 0.01\,\%  & 0.18\,\%  & 0.01\,\%  & 0.01\,\%  \\
  & Sample 2 & 11.56\,\% & 6.40\,\%  & 4.53\,\%  & 13.20\,\% & 7.44\,\%  & 4.69\,\%  \\
  & Sample 3 & 14.03\,\% & 6.84\,\%  & 6.71\,\%  & 14.55\,\% & 11.37\,\% & 8.58\,\%  \\
\multirow{3}{*}{5\,\% noise}
  & Sample 1 & 0.36\,\%  & 0.22\,\%  & 0.003\,\% & 0.11\,\%  & 0.15\,\%  & 0.003\,\% \\
  & Sample 2 & 15.29\,\% & 11.24\,\% & 8.12\,\%  & 24.44\,\% & 11.96\,\% & 8.61\,\%  \\
  & Sample 3 & 25.00\,\% & 9.20\,\%  & 8.14\,\%  & 19.03\,\% & 13.61\,\% & 10.58\,\% \\
\bottomrule
\end{tabular}
\end{center}

\vspace{1em}

% ---------- Prose: closing of 1% noise paragraph ----------
\noindent
with a 1\,\% noise level. Furthermore, we provided the values of
the regularization factors for each case depicted in Figs.~8--16
in Table~2.

\vspace{0.8em}

% ---------- Case 2 subsection heading (italic, matching paper style) ----------
\noindent\textit{Case 2: A layered ring structure}

\vspace{0.5em}

\noindent
In the second numerical example, we focused on a layered ring
structure, as illustrated in Fig.~17. To discretize the problem
domain, we employed 7610 triangular elements. The deformation of
the ring structure was induced by applying inner pressures ranging
from 10\,mmHg to 50\,mmHg, and we were only able to ``measure'' the
displacement on the inner wall. The shear modulus values of the
inner and outer layers were 0.129\,MPa and 0.386\,MPa [37],
respectively. Our approach for the inverse problem relied on
utilizing the measured displacement on the inner wall to determine
the shear modulus distribution. Initially, we provided an initial
guess of the shear modulus distribution, as shown in Fig.~18(a),
where four layers were predefined. The shear modulus values, from
the outermost layer to the innermost layer, ranged from 0.01\,MPa
to 0.04\,MPa. During the inverse problem-solving process, we
optimized both the shape of the interface between neighboring
layers and the shear modulus value of each layer. In the initial
assumption, the ring structure is considered to consist of 4
layers. For the explicit inverse method, each interface is
represented by 14 control points, leading to a total of 56
optimization variables, which includes 4 shear moduli. On the
other hand, the implicit inverse method involves a much larger
number of optimization variables, specifically 4030 variables in
total. Fig.~18 displays the variation of the reconstructed shear
modulus distribution using the explicit inverse method for the
noise free case. Notably, the reconstruction achieved through the
explicit method closely resembled the ground truth. However, the
same cannot be said for the implicit inverse method, where the
reconstructed shear modulus distribution performed significantly
worse than the explicit inverse method even for the noise free
data (see Fig.~19(d)-(f)).

% [page ends here — continues on page 12]

\newpage

% =========== Page 12: wei_page12.tex ===========

% ---------- Figure 12 placeholder ----------
\noindent
\begin{center}
\fbox{\parbox{0.92\textwidth}{\centering\vspace{1em}
\textbf{[Figure 12 --- placeholder]}\\[0.8em]
Six-panel side-by-side comparison of the \textbf{Explicit Inverse
Method} (left column) vs.\ \textbf{Implicit Inverse Method} (right
column) on \textbf{Sample 2} (the two-layer target: background
$\mu^0 \approx 3$\,MPa, top layer $\mu^1 \approx 10$\,MPa),
with \textbf{1\,\% noise} added to the measured surface
displacements, using 3, 6, and 9 displacement datasets.\\[0.4em]
\textbf{Explicit (left column)}: bilayer structure recovered with
a relatively clean (though slightly wavy) interface. Colorbar
ranges:
(a) 3 fields, SM $\in [1.00, 11.33]$\,MPa;
(b) 6 fields, SM $\in [1.00, 11.43]$\,MPa;
(c) 9 fields, SM $\in [1.00, 11.34]$\,MPa.\\[0.3em]
\textbf{Implicit (right column)}: bilayer structure visible but
with significant artifacts, especially in panel (d) where a large
green/blue patch intrudes into the red top layer from the left
edge. Colorbar ranges:
(d) 3 fields, SM $\in [1.00, 14.60]$\,MPa --- a \textbf{46\,\%
overshoot} above the nominal 10\,MPa top layer, matching the
paper's claim that the implicit method's average relative error
``could exceed 40\,\%'' under 1\,\% noise;
(e) 6 fields, SM $\in [1.00, 10.60]$\,MPa;
(f) 9 fields, SM $\in [1.00, 10.71]$\,MPa.\\[0.3em]
\textbf{Takeaway}: the noise amplifies the implicit method's
difficulty with interface recovery dramatically in the 3-field
case; 6 and 9 fields bring it back toward the nominal value but
with visibly rougher interfaces than the explicit method.
\vspace{1em}}}
\end{center}
\vspace{0.3em}
\begin{center}
\textbf{Fig. 12.} The reconstructed results with varying 3, 6 and
9 surface displacement datasets corresponding to Sample 2 (1\,\%
noise is introduced into the measured data, Unit: MPa).
\end{center}

\vspace{0.8em}

% ---------- Figure 13 placeholder ----------
\noindent
\begin{center}
\fbox{\parbox{0.92\textwidth}{\centering\vspace{1em}
\textbf{[Figure 13 --- placeholder]}\\[0.8em]
Six-panel side-by-side comparison of the \textbf{Explicit Inverse
Method} (left column) vs.\ \textbf{Implicit Inverse Method} (right
column) on \textbf{Sample 3} (the three-layer target: blue bottom
layer $\mu^0 \approx 3$\,MPa, green middle layer
$\mu^1 \approx 6$\,MPa, red top layer $\mu^2 \approx 10$\,MPa),
with \textbf{1\,\% noise} added to the measured surface
displacements, using 3, 6, and 9 displacement datasets.\\[0.4em]
\textbf{Explicit (left column)}: three-layer structure recovered
with wavy but distinct interfaces between the layers. Colorbar
ranges:
(a) 3 fields, SM $\in [1.01, 10.63]$\,MPa;
(b) 6 fields, SM $\in [1.00, 10.25]$\,MPa;
(c) 9 fields, SM $\in [1.00, 10.38]$\,MPa.\\[0.3em]
\textbf{Implicit (right column)}: three-layer structure noticeably
noisier, with fuzzier interfaces and small artifact patches.
Colorbar ranges:
(d) 3 fields, SM $\in [1.00, 11.78]$\,MPa;
(e) 6 fields, SM $\in [1.00, 11.58]$\,MPa;
(f) 9 fields, SM $\in [1.00, 9.96]$\,MPa --- notably, the 9-field
implicit result \textbf{undershoots} the nominal 10\,MPa target,
unusual relative to most other panels in the paper which tend to
overshoot.\\[0.3em]
\textbf{Takeaway}: as with Sample 2 under noise, the explicit
method produces cleaner interfaces; but the quantitative
colorbar-max values on Sample 3 are closer between explicit and
implicit than on Sample 2 because the three-layer structure
constrains both methods more tightly than the two-layer structure.
\vspace{1em}}}
\end{center}
\vspace{0.3em}
\begin{center}
\textbf{Fig. 13.} The reconstructed results with varying 3, 6 and
9 surface displacement datasets corresponding to Sample 3 (1\,\%
noise is introduced into the measured data, Unit: MPa).
\end{center}

\vspace{1em}

% ---------- Prose paragraph (from pdftotext, soft-hyphen cleaned) ----------
\noindent
When 1\,\% noise was introduced, the explicit inverse method
continued to perform admirably, producing high-quality
reconstructed results (see Fig.~20). However, the performance of
the implicit inverse method continued to be inferior to that of
the explicit inverse method. The average relative error for the
implicit inverse method could exceed 40\,\%, indicating substantial
inaccuracies in the reconstructed results. On the other hand, the
average relative error for the explicit inverse method remained
relatively low, around 5\,\%, signifying its superior ability to
handle the noise and produce more accurate reconstructions (see
Table~3). The values of regularization factors used in each case
from Fig.~19 and Fig.~20 are listed in Table~4.

% [page ends here — continues on page 13]

\newpage

% =========== Page 13: wei_page13.tex ===========

% ---------- Figure 14 placeholder ----------
\noindent
\begin{center}
\fbox{\parbox{0.92\textwidth}{\centering\vspace{1em}
\textbf{[Figure 14 --- placeholder]}\\[0.8em]
Six-panel side-by-side comparison of the \textbf{Explicit Inverse
Method} (left column) vs.\ \textbf{Implicit Inverse Method} (right
column) on \textbf{Sample 1} (the homogeneous target, true
SM $= 3.00$\,MPa), with \textbf{5\,\% noise} added to the measured
surface displacements, using 3, 6, and 9 displacement datasets.\\[0.4em]
\textbf{Explicit (left column)}: the higher noise level corrupts
the homogeneous recovery. Colorbar ranges:
(a) 3 fields, SM $\in [2.9641, 3.0096]$\,MPa --- visible spurious
\textbf{bilayer artifact} with a blue ``valley'' cutting through
an otherwise red field (the homogeneous target should be uniform);
(b) 6 fields, SM $\in [3.0014, 3.0099]$\,MPa --- wavy spurious
bilayer pattern, red top over blue bottom;
(c) 9 fields, SM $\in [3.0001, 3.0001]$\,MPa --- uniform blue, near
perfect. The 9-field case recovers cleanly while 3-field and
6-field hallucinate structure that isn't there.\\[0.3em]
\textbf{Implicit (right column)}:
(d) 3 fields, SM $\in [2.7156, 3.0220]$\,MPa --- mostly uniform red
with a small localized blue blob;
(e) 6 fields, SM $\in [3.0027, 3.0060]$\,MPa --- visible vertical
gradient/striping;
(f) 9 fields, SM $\in [2.9999, 3.0002]$\,MPa --- almost uniform with
faint artifact pattern.\\[0.3em]
\textbf{Takeaway}: under 5\,\% noise on a homogeneous target, the
explicit method's 3-field and 6-field reconstructions show
\emph{spurious bilayer patterns} that don't exist in the ground
truth. This is a failure mode worth flagging: the method can
invent structure under noise when very few displacement datasets
are available. Nine fields restores clean recovery.
\vspace{1em}}}
\end{center}
\vspace{0.3em}
\begin{center}
\textbf{Fig. 14.} The reconstructed results with varying 3, 6 and
9 surface displacement datasets corresponding to Sample 1 (5\,\%
noise is introduced into the measured data, Unit: MPa).
\end{center}

\vspace{0.8em}

% ---------- Figure 15 placeholder ----------
\noindent
\begin{center}
\fbox{\parbox{0.92\textwidth}{\centering\vspace{1em}
\textbf{[Figure 15 --- placeholder]}\\[0.8em]
Six-panel side-by-side comparison of the \textbf{Explicit Inverse
Method} (left column) vs.\ \textbf{Implicit Inverse Method} (right
column) on \textbf{Sample 2} (the two-layer target: background
$\mu^0 \approx 3$\,MPa, top layer $\mu^1 \approx 10$\,MPa),
with \textbf{5\,\% noise} added, using 3, 6, and 9 displacement
datasets.\\[0.4em]
\textbf{Explicit (left column)}: bilayer structure visible in all
three panels. Colorbar ranges:
(a) 3 fields, SM $\in [1.00, 8.91]$\,MPa --- notably, the colorbar
max is \textbf{11\,\% below the nominal 10\,MPa} top layer. This
is an \emph{undershoot}, opposite to the usual overshoot failure
mode. The top layer is not fully recovered with only 3 fields
under 5\,\% noise;
(b) 6 fields, SM $\in [1.00, 10.87]$\,MPa --- slight overshoot,
wavy interface;
(c) 9 fields, SM $\in [1.00, 10.16]$\,MPa --- near-nominal with
slight overshoot.\\[0.3em]
\textbf{Implicit (right column)}:
(d) 3 fields, SM $\in [1.00, 14.20]$\,MPa --- \textbf{42\,\%
overshoot} above the 10\,MPa nominal, visible with significant
green/yellow patches intruding into the top layer;
(e) 6 fields, SM $\in [1.00, 11.42]$\,MPa;
(f) 9 fields, SM $\in [1.00, 11.43]$\,MPa --- surprising that 9
fields is not clearly better than 6 fields for the implicit
method on this sample.\\[0.3em]
\textbf{Takeaway}: the 3-field explicit case \emph{undershoots}
while the 3-field implicit case \emph{overshoots} by 42\,\% ---
the two methods fail in opposite directions under low-data +
5\,\% noise. More displacement fields close the gap for explicit
but not cleanly for implicit.
\vspace{1em}}}
\end{center}
\vspace{0.3em}
\begin{center}
\textbf{Fig. 15.} The reconstructed results with varying 3, 6 and
9 surface displacement datasets corresponding to Sample 2 (5\,\%
noise is introduced into the measured data, Unit: MPa).
\end{center}

% [page ends here — continues on page 14]

\newpage

% =========== Page 14: wei_page14.tex ===========

% ---------- Figure 16 placeholder ----------
\noindent
\begin{center}
\fbox{\parbox{0.92\textwidth}{\centering\vspace{1em}
\textbf{[Figure 16 --- placeholder]}\\[0.8em]
Six-panel side-by-side comparison of the \textbf{Explicit Inverse
Method} (left column) vs.\ \textbf{Implicit Inverse Method} (right
column) on \textbf{Sample 3} (the three-layer target: blue bottom
$\mu^0 \approx 3$\,MPa, green middle $\mu^1 \approx 6$\,MPa, red
top $\mu^2 \approx 10$\,MPa), with \textbf{5\,\% noise} added,
using 3, 6, and 9 displacement datasets.\\[0.4em]
\textbf{Explicit (left column)}:
(a) 3 fields, SM $\in [1.00, 8.46]$\,MPa --- undershoots the
10\,MPa nominal top layer, matching the 3-field undershoot pattern
seen in Fig.~15 for Sample 2;
(b) 6 fields, SM $\in [1.00, 10.28]$\,MPa;
(c) 9 fields, SM $\in [1.00, 10.34]$\,MPa.
All three panels show wavy but distinguishable three-layer
structure.\\[0.3em]
\textbf{Implicit (right column)}:
(d) 3 fields, SM $\in [1.00, 15.49]$\,MPa --- 55\,\% overshoot,
significant green/yellow patches corrupting the top red layer;
(e) 6 fields, SM $\in [1.00, 15.98]$\,MPa --- \textbf{worse than
3 fields} on the max-SM axis, a \emph{non-monotone} degradation
that the paper's prose does not call out;
(f) 9 fields, SM $\in [1.00, 10.17]$\,MPa --- snaps back close to
nominal, cleaner interfaces.\\[0.3em]
\textbf{Takeaway}: the implicit method's colorbar max on Sample 3
under 5\,\% noise is \emph{not monotone} in the number of
displacement datasets ($15.49 \to 15.98 \to 10.17$ across 3, 6, 9
fields). This is unusual in the paper's overall story; worth
noting as a caveat to any ``more data always helps'' interpretation.
\vspace{1em}}}
\end{center}
\vspace{0.3em}
\begin{center}
\textbf{Fig. 16.} The reconstructed results with varying 3, 6 and
9 surface displacement datasets corresponding to Sample 3 (5\,\%
noise is introduced into the measured data, Unit: MPa).
\end{center}

\vspace{1em}

% ---------- Table 2 ----------
\noindent\textbf{Table 2}\\[0.2em]
{\small The values of the regularization factors utilized in
Fig.~8, Fig.~9, Fig.~10, Fig.~11, Fig.~12, Fig.~13, Fig.~14,
Fig.~15, Fig.~16.}
\vspace{0.3em}

\begin{center}
\footnotesize
\begin{tabular}{@{}ll ccc ccc@{}}
\toprule
\multicolumn{2}{l}{Inverse method}
  & \multicolumn{3}{c}{Explicit inverse method}
  & \multicolumn{3}{c}{Implicit inverse method} \\
\cmidrule(lr){3-5}\cmidrule(lr){6-8}
\multicolumn{2}{l}{Number of surface displacement datasets}
  & 3 & 6 & 9 & 3 & 6 & 9 \\
\midrule
\multirow{3}{*}{Sample 1}
  & no noise  & 0.00            & 0.00            & $5\times 10^{-27}$ & $2\times 10^{-14}$ & $5\times 10^{-14}$ & $1\times 10^{-13}$ \\
  & 1\,\% noise & $1\times 10^{-11}$ & $1\times 10^{-11}$ & $1\times 10^{-18}$ & $2\times 10^{-11}$ & $6\times 10^{-11}$ & $6\times 10^{-10}$ \\
  & 5\,\% noise & $5\times 10^{-12}$ & $1\times 10^{-10}$ & $1\times 10^{-10}$ & $1\times 10^{-10}$ & $6\times 10^{-10}$ & $6\times 10^{-9}$  \\
\midrule
\multirow{3}{*}{Sample 2}
  & no noise  & 0.00            & 0.00            & 0.00            & $1\times 10^{-14}$ & $1\times 10^{-14}$ & $1\times 10^{-14}$ \\
  & 1\,\% noise & $1\times 10^{-13}$ & $1\times 10^{-13}$ & $1\times 10^{-13}$ & $1\times 10^{-12}$ & $1\times 10^{-11}$ & $1\times 10^{-11}$ \\
  & 5\,\% noise & $1\times 10^{-10}$ & $1\times 10^{-10}$ & $1\times 10^{-10}$ & $1\times 10^{-11}$ & $1\times 10^{-10}$ & $1\times 10^{-10}$ \\
\midrule
\multirow{3}{*}{Sample 3}
  & no noise  & 0.00            & 0.00            & 0.00            & $1\times 10^{-14}$ & $1\times 10^{-14}$ & $1\times 10^{-14}$ \\
  & 1\,\% noise & 0.00            & 0.00            & 0.00            & $1\times 10^{-12}$ & $1\times 10^{-12}$ & $1\times 10^{-11}$ \\
  & 5\,\% noise & $1\times 10^{-11}$ & $1\times 10^{-11}$ & $1\times 10^{-11}$ & $1\times 10^{-11}$ & $1\times 10^{-11}$ & $1\times 10^{-10}$ \\
\bottomrule
\end{tabular}
\end{center}

\vspace{1em}

% ---------- Figure 17 placeholder ----------
\noindent
\begin{center}
\fbox{\parbox{0.80\textwidth}{\centering\vspace{1em}
\textbf{[Figure 17 --- placeholder]}\\[0.8em]
Target shear modulus distribution for the \textbf{Case 2} bilayer
ring structure: a circular ring with an outer thicker red layer
($\mu_\text{outer} = 0.386$\,MPa) surrounding an inner thinner
blue layer ($\mu_\text{inner} = 0.129$\,MPa). Black arrows point
radially inward around the perimeter, labeled ``displacement
measurement'', indicating that the inverse problem uses only the
displacement measured on the inner wall. Colorbar ticks at
$\{0.129, 0.200, 0.267, 0.333, 0.386\}$\,MPa (jet colormap,
matching the values $[37]$ for physiological soft tissue from the
paper's bibliography). The stiffness ratio outer/inner
$\approx 3.0$ is realistic for arterial wall tissue.
\vspace{1em}}}
\end{center}
\vspace{0.3em}
\begin{center}
\textbf{Fig. 17.} Target shear modulus distribution of a bilayer
ring structure (Unit: MPa).
\end{center}

% [page ends here — continues on page 15]

\newpage

% =========== Page 15: wei_page15.tex ===========

% ---------- Figure 18 placeholder ----------
\noindent
\begin{center}
\fbox{\parbox{0.92\textwidth}{\centering\vspace{1em}
\textbf{[Figure 18 --- placeholder]}\\[0.8em]
Six-panel optimization trajectory for the \textbf{Case 2 bilayer
ring structure} (target: outer 0.386\,MPa, inner 0.129\,MPa) using
the explicit inverse scheme with 3 surface displacement datasets
(no noise), starting from a 4-layer initial guess. Each panel is
a circular ring with a jet colorbar.\\[0.4em]
\begin{tabular}{@{}ll@{}}
\textbf{(a) Initial guess}  & 4-layer ring, SM $\in [0.01, 0.04]$\,MPa \\
\textbf{(b) Iteration 20}   & SM $\in [0.145, 0.438]$\,MPa (overshoots outer layer) \\
\textbf{(c) Iteration 50}   & SM $\in [0.131, 0.397]$\,MPa, wavy interface emerging \\
\textbf{(d) Iteration 100}  & SM $\in [0.129, 0.393]$\,MPa \\
\textbf{(e) Iteration 1000} & SM $\in [0.131, 0.387]$\,MPa (near target) \\
\textbf{(f) Iteration 1322 (final)} & SM $\in [0.131, 0.384]$\,MPa
\end{tabular}\\[0.4em]
The outer-layer colorbar max converges to 0.386\,MPa
\emph{from above}: the trajectory overshoots early
($0.438 \to 0.397 \to 0.393 \to 0.387 \to 0.384$) and settles
slightly under the nominal 0.386 target. The ring interface
becomes increasingly wavy as iteration proceeds but the overall
bilayer structure is recovered faithfully.
\vspace{1em}}}
\end{center}
\vspace{0.3em}
\begin{center}
\textbf{Fig. 18.} The optimization process with four layers
initially utilizing the explicit inverse scheme with 3 surface
datasets (Without noise, Unit: MPa).
\end{center}

\vspace{0.8em}

% ---------- Figure 19 placeholder ----------
\noindent
\begin{center}
\fbox{\parbox{0.92\textwidth}{\centering\vspace{1em}
\textbf{[Figure 19 --- placeholder]}\\[0.8em]
Six-panel comparison of the \textbf{Explicit Inverse Method} (top
row) vs.\ \textbf{Implicit Inverse Method} (bottom row) on the
bilayer ring structure (\textbf{no noise}), with 3, 6, and 9
surface displacement datasets.\\[0.4em]
\textbf{Explicit (top row)}: all three reconstructions are
visually near-identical to the Fig.~17 target --- clean bilayer
ring with wavy inner interface. Colorbar ranges:
(a) 3 fields, SM $\in [0.131, 0.384]$\,MPa;
(b) 6 fields, SM $\in [0.131, 0.384]$\,MPa;
(c) 9 fields, SM $\in [0.131, 0.384]$\,MPa.
The explicit method reaches the same quality with 3 fields as
with 9 --- consistent with its low parameter count
(56 variables per the paper's earlier claim).\\[0.4em]
\textbf{Implicit (bottom row)}: reconstructions are visibly
\textbf{dominated by green/yellow artifacts} with no clear bilayer
structure recovered. Colorbar ranges show a \emph{monotonically
increasing max} as more datasets are added:
(d) 3 fields, SM $\in [0.010, 0.478]$\,MPa;
(e) 6 fields, SM $\in [0.010, 0.632]$\,MPa;
(f) 9 fields, SM $\in [0.010, 0.813]$\,MPa.
More data makes the implicit result \emph{worse}, not better ---
the 9-field max overshoots the 0.386 nominal outer-layer value by
\textbf{110\,\%}.\\[0.3em]
\textbf{Takeaway}: this is the most dramatic failure mode of the
implicit method in the entire paper. On the ring structure, with
4030 optimization variables to the explicit method's 56, the
implicit method cannot find the bilayer structure at all ---
increasing displacement data only makes it diverge more. This
suggests the implicit method's high dimensionality is
\emph{actively harmful} on this problem, not just inefficient.
\vspace{1em}}}
\end{center}
\vspace{0.3em}
\begin{center}
\textbf{Fig. 19.} The reconstructed results with varying 3, 6 and
9 surface displacement datasets (Without noise, Unit: MPa).
\end{center}

\vspace{1em}

% ---------- Case 3 subsection heading (italic, matching paper style) ----------
\noindent\textit{Case 3: A composite with inclusions embedded}

\vspace{0.5em}

\noindent
In the last example, we replicated a scenario where one or two
tumors were embedded within soft tissues with a stiffness ratio
of 5:1, as illustrated in Fig.~21. To simulate the forward problem,
we applied an ``indentation'' force on the top surface of a square
domain, discretized using 7200 triangular elements. The center
point was fixed and other points where on the bottom edge are
constrained in the vertical direction, and we only utilized the
displacements on the top surface to tackle the inverse problem.
To initiate the inverse problem, we employed four inclusions as
the initial guess (refer to Fig.~22(a)). These inclusions were
predefined in the problem

% [page ends here — continues on page 16]

\newpage

% =========== Page 16: wei_page16.tex ===========

% ---------- Figure 20 placeholder ----------
\noindent
\begin{center}
\fbox{\parbox{0.92\textwidth}{\centering\vspace{1em}
\textbf{[Figure 20 --- placeholder]}\\[0.8em]
Six-panel comparison of \textbf{Explicit Inverse Method} (top row)
vs.\ \textbf{Implicit Inverse Method} (bottom row) on the
\textbf{Case 2 bilayer ring structure} with \textbf{1\,\% noise}
added, using 3, 6, and 9 surface displacement datasets.\\[0.4em]
\textbf{Explicit (top row)}: near-identical bilayer ring
reconstructions, visually indistinguishable from the Fig.~17
target. Colorbar ranges:
(a) 3 fields, SM $\in [0.129, 0.385]$\,MPa;
(b) 6 fields, SM $\in [0.132, 0.384]$\,MPa;
(c) 9 fields, SM $\in [0.130, 0.386]$\,MPa.
The explicit method tolerates 1\,\% noise essentially perfectly
on the ring problem, with all three panels landing within
$\pm 0.002$\,MPa of the 0.386 nominal outer layer.\\[0.4em]
\textbf{Implicit (bottom row)}: severe overshoot pattern.
Colorbar ranges:
(d) 3 fields, SM $\in [0.010, 1.251]$\,MPa --- \textbf{224\,\%
overshoot} above the 0.386 nominal outer layer; visually dominated
by green/yellow artifacts;
(e) 6 fields, SM $\in [0.010, 0.874]$\,MPa --- 126\,\% overshoot;
(f) 9 fields, SM $\in [0.010, 0.975]$\,MPa --- 152\,\% overshoot.
The implicit method's max is also non-monotone in this figure
($1.251 \to 0.874 \to 0.975$), continuing the pattern from Fig.~16
and Fig.~19.\\[0.3em]
\textbf{Takeaway}: as in the noise-free case (Fig.~19), the
implicit method is fundamentally failing on the ring problem. The
1\,\% noise addition does not change the qualitative picture ---
it is already broken without noise. See Table~3 for the
quantitative confirmation (implicit $\approx 41$\,\% vs.\ explicit
$\approx 5$\,\% average relative error).
\vspace{1em}}}
\end{center}
\vspace{0.3em}
\begin{center}
\textbf{Fig. 20.} The reconstructed results with varying 3, 6 and
9 surface displacement datasets (1\,\% noise is introduced,
Unit: MPa).
\end{center}

\vspace{1em}

% ---------- Table 3 ----------
\noindent\textbf{Table 3}\\[0.2em]
{\small Comparison of the average relative error in the shear
modulus distribution obtained by implicit inverse and explicit
inverse methods for the ring structure (1\,\% noise is introduced).}
\vspace{0.3em}

\begin{center}
\small
\begin{tabular}{@{}l ccc ccc@{}}
\toprule
Inverse method
  & \multicolumn{3}{c}{Explicit inverse method}
  & \multicolumn{3}{c}{Implicit inverse method} \\
\cmidrule(lr){2-4}\cmidrule(lr){5-7}
Number of surface displacement datasets
  & 3 & 6 & 9 & 3 & 6 & 9 \\
\midrule
Average relative error
  & 4.73\,\% & 5.38\,\% & 4.56\,\%
  & 42.41\,\% & 40.99\,\% & 41.23\,\% \\
\bottomrule
\end{tabular}
\end{center}

\vspace{1em}

% ---------- Table 4 ----------
\noindent\textbf{Table 4}\\[0.2em]
{\small The values of regularization factors used in Figs.~19 and 20.}
\vspace{0.3em}

\begin{center}
\small
\begin{tabular}{@{}ll ccc ccc@{}}
\toprule
\multicolumn{2}{l}{Inverse method}
  & \multicolumn{3}{c}{Explicit inverse method}
  & \multicolumn{3}{c}{Implicit inverse method} \\
\cmidrule(lr){3-5}\cmidrule(lr){6-8}
\multicolumn{2}{l}{Number of surface displacement datasets}
  & 3 & 6 & 9 & 3 & 6 & 9 \\
\midrule
\multirow{2}{*}{Ring structure}
  & no noise    & 0.00 & 0.00 & 0.00 & $1\times 10^{-12}$ & $1\times 10^{-12}$ & $1\times 10^{-12}$ \\
  & 1\,\% noise & 0.00 & 0.00 & 0.00 & $1\times 10^{-12}$ & $1\times 10^{-12}$ & $1\times 10^{-12}$ \\
\bottomrule
\end{tabular}
\end{center}

\vspace{1em}

% ---------- Figure 21 placeholder ----------
\noindent
\begin{center}
\fbox{\parbox{0.85\textwidth}{\centering\vspace{1em}
\textbf{[Figure 21 --- placeholder]}\\[0.8em]
\textbf{Target shear modulus distribution} for the \textbf{Case 3
composite with embedded inclusions}. Two sub-panels:\\[0.3em]
\textbf{(a) Sample 1}: a square blue background (soft tissue,
$\mu_\text{background} = 1.00$\,MPa) with a single red circular
inclusion at the center (tumor, $\mu_\text{inclusion} = 5.00$\,MPa).
Stiffness ratio 5:1, matching the Case~3 setup described on
page~140.\\[0.3em]
\textbf{(b) Sample 2}: same square background with \textbf{two}
red circular inclusions side-by-side in the upper half of the
domain.\\[0.3em]
Both panels share boundary conditions: the top edge has an array
of small downward-pointing arrows labeled ``displacement
measurement'' (these are the indentation-force application points
and where surface displacements are recorded for the inverse
problem); the bottom edge is marked with small circles/triangles
indicating fixed/vertical-direction constraints. Colorbar ticks
at $\{1.00, 2.00, 3.00, 4.00, 5.00\}$\,MPa, jet colormap.
\vspace{1em}}}
\end{center}
\vspace{0.3em}
\begin{center}
\textbf{Fig. 21.} Target shear modulus distribution of composite
with inclusions embedded (Unit: MPa).
\end{center}

% [page ends here — continues on page 17]

\newpage

% =========== Page 17: wei_page17.tex ===========

% ---------- Figure 22 placeholder ----------
\noindent
\begin{center}
\fbox{\parbox{0.92\textwidth}{\centering\vspace{1em}
\textbf{[Figure 22 --- placeholder]}\\[0.8em]
Six-panel optimization trajectory for the \textbf{Case 3 Sample 1}
(one embedded inclusion, target SM $\approx 5$\,MPa against a
1\,MPa soft-tissue background), using the explicit inverse scheme
with 3 surface displacement datasets (no noise). The initial
guess has \textbf{four} inclusions predefined at the four corners
of the square domain; the optimization merges them into a single
central inclusion as it converges.\\[0.4em]
\begin{tabular}{@{}ll@{}}
\textbf{(a) Initial guess}    & 4 green circles at corners, SM $\in [1.00, 1.00]$\,MPa (uniform) \\
\textbf{(b) Iteration 5}      & 4 merging blobs, SM $\in [1.00, 3.20]$\,MPa \\
\textbf{(c) Iteration 10}     & 3 blobs (center two merging), SM $\in [1.00, 3.23]$\,MPa \\
\textbf{(d) Iteration 15}     & 1--2 blobs (red central), SM $\in [1.00, 4.26]$\,MPa \\
\textbf{(e) Iteration 275}    & 1 red inclusion, SM $\in [2.10, 5.32]$\,MPa \\
\textbf{(f) Iteration 328 (final)} & 1 red inclusion, SM $\in [1.00, 5.10]$\,MPa
\end{tabular}\\[0.4em]
\textbf{Takeaway}: the 4-inclusion initialization does NOT cause
the algorithm to find 4 tumors. Instead, the four components
merge into the single ground-truth tumor by iteration 15. This
is the opposite of the usual over-parameterization concern:
redundant initial components collapse to the correct solution
rather than forcing a spurious multi-component reconstruction.
The explicit parameterization has 69 total variables
(4 inclusions $\times$ 16 geometric parameters + 4 shear moduli +
1 background shear modulus).
\vspace{1em}}}
\end{center}
\vspace{0.3em}
\begin{center}
\textbf{Fig. 22.} The optimization process with four inclusions
initially utilizing the explicit inverse scheme for a solid with
inclusion embedded when 3 surface displacement datasets are
utilized (Without noise, Unit: MPa).
\end{center}

\vspace{0.8em}

% ---------- Prose paragraph (continuation from page 16) ----------
\noindent
domain. Each inclusion was represented by 16 geometric parameters
and 1 material property, which is the shear modulus. Consequently,
the total number of optimization variables reached 69, inclusive
of the shear modulus value of the background.

During the optimization process, the four predefined components
gradually merged into one single inclusion. As convergence was
reached, the shear modulus of this inclusion approached the
target values effectively, demonstrating the success of the
inverse problem-solving method. The reconstructed results, as
shown in Fig.~22, illustrated the accurate estimation of the
shear modulus for the inclusion after the convergence of the
algorithm. Indeed, the explicit inverse method demonstrated a
notably faster convergence rate compared to the implicit inverse
method. Specifically, the explicit inverse method achieved
convergence in just 328 iterations,

\vspace{0.8em}

% ---------- Figure 23 placeholder ----------
\noindent
\begin{center}
\fbox{\parbox{0.92\textwidth}{\centering\vspace{1em}
\textbf{[Figure 23 --- placeholder]}\\[0.8em]
Six-panel comparison of the \textbf{Explicit Inverse Method} (top
row) vs.\ \textbf{Implicit Inverse Method} (bottom row) on
\textbf{Case 3 Sample 1} (one embedded inclusion, target
$\approx 5$\,MPa), with \textbf{no noise}, using 3, 6, and 9
surface displacement datasets.\\[0.4em]
\textbf{Explicit (top row)}:
(a) 3 fields, SM $\in [1.00, 5.10]$\,MPa;
(b) 6 fields, SM $\in [1.00, 5.06]$\,MPa;
(c) 9 fields, SM $\in [1.00, 4.97]$\,MPa.
All three panels recover a clean single red inclusion on a blue
background, near the 5\,MPa nominal target. The colorbar max
drifts slightly downward as more fields are added
($5.10 \to 5.06 \to 4.97$).\\[0.4em]
\textbf{Implicit (bottom row)}:
(d) 3 fields, SM $\in [1.00, 4.47]$\,MPa;
(e) 6 fields, SM $\in [1.00, 4.58]$\,MPa;
(f) 9 fields, SM $\in [1.00, 4.79]$\,MPa.
\textbf{All three panels undershoot the 5\,MPa nominal target},
with the 3-field case falling short by 10.6\,\%. The max values
are non-monotone in dataset count
($4.47 \to 4.58 \to 4.79$) but monotonically approach the target
as more fields are added --- the opposite direction from the
explicit method's slight drift.\\[0.3em]
\textbf{Takeaway}: on the one-inclusion case without noise, both
methods produce visually similar reconstructions, but the
explicit method hits the target nominal from above
(slight overshoot) while the implicit method hits it from below
(persistent undershoot). More data helps the implicit method
approach nominal but doesn't quite reach it.
\vspace{1em}}}
\end{center}
\vspace{0.3em}
\begin{center}
\textbf{Fig. 23.} The reconstructed results with varying 3, 6 and
9 surface displacement datasets corresponding to Sample 1
(Without noise, Unit: MPa).
\end{center}

% [page ends here — continues on page 18]

\newpage

% =========== Page 18: wei_page18.tex ===========

\noindent
while the implicit inverse method required a significantly higher
number of iterations, reaching convergence after 4354 iterations.

In the absence of noise in the datasets, both the explicit and
implicit inverse methods exhibited excellent performance in
yielding shear modulus reconstructions, as demonstrated in
Fig.~23. However, when 1\,\% noise was introduced to the
datasets, the explicit inverse method outperformed the implicit
inverse method (see Fig.~24). For the explicit inverse method,
the average relative error in the reconstructed shear modulus
values was approximately 2\,\%, indicating a relatively accurate
estimation despite the presence of noise (see Table~5). On the
other hand, the implicit inverse method exhibited a higher
average error of more than 3\,\%, signifying a less precise
reconstruction in the presence of noise. Overall, the explicit
inverse method proved to be more robust and capable of handling
noisy datasets, leading to more accurate shear modulus
reconstructions compared to the implicit inverse method.
Furthermore, it is noteworthy that the reconstructed results
maintain good quality even when the noise level is increased to
5\,\% (see Fig.~25).

In the case where two inclusions were embedded in the background
(as depicted in Fig.~21(b)). As shown in Fig.~26, the results
obtained using the implicit inverse method showed a consistent
underestimation of the recovered shear modulus values by at
least 20\,\% as the size of the mapped inclusions increased.
Additionally, the shape of the mapped inclusions appeared to
become more circular with the increase of surface displacement
datasets. Conversely, when employing the explicit inverse
method, the shear modulus values within the two inclusions were
overestimated by approximately 20\,\%. Furthermore, the shape
of the recovered inclusions exhibited a more circular appearance
compared to those predicted by the implicit inverse method. The
regularization factors used in Figs.~23--26 are listed in
Table~6.

To assess the robustness of the proposed explicit inverse
method, we varied the shape of the inclusion in Sample~1, as
depicted in Fig.~27(a) and~(d). Our observations indicate that
the reconstructed results obtained from the explicit inverse
method effectively preserve the shape of the complex inclusion,
even with only three displacement datasets. In contrast, the
implicit inverse method failed to achieve this level of
accuracy.

\vspace{0.5em}

\noindent\textbf{Discussion}

\vspace{0.3em}

\noindent
The paper proposes an explicit inverse approach for
characterizing the nonhomogeneous elastic property distribution
of three typical nonlinear biological structures with surface
deformation. This approach is different from the prevalent
implicit inverse approach in that it optimizes only the
geometric parameters and shear moduli of each region, reducing
the scale of the optimization problem. By doing so, the proposed
method improves the uniqueness of the inverse solution since the
total number of optimization variables is reduced. Additionally,
the paper highlights an important advantage of the explicit
method: it only requires surface deformation data, which can be
measured with high accuracy using low-cost measurement
techniques. This makes the proposed approach more practical and
accessible for real-world applications, as it avoids the need
for complex and expensive instrumentation to gather the
full-field data.

\vspace{0.8em}

% ---------- Figure 24 placeholder ----------
\noindent
\begin{center}
\fbox{\parbox{0.92\textwidth}{\centering\vspace{1em}
\textbf{[Figure 24 --- placeholder]}\\[0.8em]
Six-panel comparison of \textbf{Explicit Inverse Method} (top
row) vs.\ \textbf{Implicit Inverse Method} (bottom row) on
\textbf{Case 3 Sample 1} (one embedded inclusion, target
$\approx 5$\,MPa), with \textbf{1\,\% noise} added, using 3, 6,
and 9 surface displacement datasets.\\[0.4em]
\textbf{Explicit (top row)}:
(a) 3 fields, SM $\in [1.00, 5.75]$\,MPa --- \textbf{15\,\%
overshoot} of the 5\,MPa nominal tumor stiffness;
(b) 6 fields, SM $\in [1.00, 5.30]$\,MPa;
(c) 9 fields, SM $\in [1.00, 5.01]$\,MPa --- near-perfect with
9 fields. The max drifts downward monotonically
($5.75 \to 5.30 \to 5.01$) toward the nominal target.\\[0.4em]
\textbf{Implicit (bottom row)}:
(d) 3 fields, SM $\in [1.00, 3.26]$\,MPa --- \textbf{35\,\%
undershoot}, a dramatically worse number than the paper's prose
claim of ``$\approx 3$\,\% average error'' suggests (the prose
error is measured over the whole domain, not the inclusion
max);
(e) 6 fields, SM $\in [1.00, 3.51]$\,MPa --- 30\,\% undershoot;
(f) 9 fields, SM $\in [1.00, 4.82]$\,MPa --- still 3.6\,\% below
nominal. Max drifts upward monotonically
($3.26 \to 3.51 \to 4.82$) toward the target but doesn't
reach it.\\[0.3em]
\textbf{Takeaway}: with 1\,\% noise, the explicit and implicit
methods fail in \emph{opposite directions} on this problem ---
explicit \emph{overshoots} the tumor stiffness, implicit
\emph{undershoots} it. Both approach the 5\,MPa target from
their respective sides as more fields are added, but the
explicit method gets closer faster. The prose's ``2\,\% vs
3\,\% average error'' framing understates the practical
significance: a 35\,\% undershoot of the tumor stiffness
(panel d) would have very different clinical implications than
a 2\,\% average-domain error suggests.
\vspace{1em}}}
\end{center}
\vspace{0.3em}
\begin{center}
\textbf{Fig. 24.} The reconstructed results with varying 3, 6
and 9 surface displacement datasets corresponding to Sample 1
(1\,\% noise is introduced into the measured data, Unit: MPa).
\end{center}

% [page ends here — continues on page 19]

\newpage

% =========== Page 19: wei_page19.tex ===========

% ---------- Table 5 ----------
\noindent\textbf{Table 5}\\[0.2em]
{\small Comparison of the average relative error in the shear
modulus distribution obtained by implicit inverse and explicit
inverse methods for the case with one inclusion embedded
(Figs.~24 and~25).}
\vspace{0.3em}

\begin{center}
\small
\begin{tabular}{@{}l ccc ccc@{}}
\toprule
Inverse method
  & \multicolumn{3}{c}{Explicit inverse method}
  & \multicolumn{3}{c}{Implicit inverse method} \\
\cmidrule(lr){2-4}\cmidrule(lr){5-7}
Number of surface displacement datasets
  & 3 & 6 & 9 & 3 & 6 & 9 \\
\midrule
1\,\% noise & 1.81\,\% & 1.50\,\% & 1.60\,\% & 4.28\,\% & 3.58\,\% & 3.41\,\% \\
5\,\% noise & 1.88\,\% & 1.86\,\% & 1.74\,\% & 6.49\,\% & 6.05\,\% & 4.26\,\% \\
\bottomrule
\end{tabular}
\end{center}

\vspace{1em}

% ---------- Figure 25 placeholder ----------
\noindent
\begin{center}
\fbox{\parbox{0.92\textwidth}{\centering\vspace{1em}
\textbf{[Figure 25 --- placeholder]}\\[0.8em]
Six-panel comparison of \textbf{Explicit Inverse Method} (top
row) vs.\ \textbf{Implicit Inverse Method} (bottom row) on
Sample 1 with \textbf{5\,\% noise} added, using 3, 6, and 9
displacement datasets.\\[0.4em]
\textbf{Explicit (top row)}:
(a) 3 fields, SM $\in [1.00, 5.83]$\,MPa --- 17\,\% overshoot;
(b) 6 fields, SM $\in [1.00, 6.09]$\,MPa --- \textbf{22\,\%
overshoot, WORSE than 3-field};
(c) 9 fields, SM $\in [1.00, 4.66]$\,MPa --- now 7\,\%
\emph{undershoot}. The max is non-monotone
($5.83 \to 6.09 \to 4.66$) and crosses the nominal 5\,MPa target
between 6 and 9 fields. This contradicts the intuitive ``more
data $\Rightarrow$ better'' story.\\[0.4em]
\textbf{Implicit (bottom row)}:
(d) 3 fields, SM $\in [1.00, 2.90]$\,MPa --- \textbf{42\,\%
undershoot};
(e) 6 fields, SM $\in [1.00, 3.21]$\,MPa;
(f) 9 fields, SM $\in [1.00, 3.40]$\,MPa --- still 32\,\%
undershoot.
The implicit method never approaches the 5\,MPa target under
5\,\% noise; even 9 fields leaves it severely below
nominal.\\[0.3em]
\textbf{Table 5 vs. this figure --- an important disconnect}:
Table 5 reports the explicit-9-fields-5\%-noise case as
\emph{1.74\,\% average relative error}, but panel (c) shows the
tumor-stiffness max at \emph{4.66 instead of 5.00 MPa}. Both
numbers are correct --- Table 5 averages over the whole 2D
domain (which is mostly 1\,MPa background where any method is
trivially correct), while the figure shows the peak inside the
tumor itself. A clinician looking at the ``$\approx 2\,\%$
error'' claim would expect a $\pm 0.1$\,MPa estimate of the
tumor; the figure shows a $-7\,\%$ bias in the quantity that
actually matters. The prose should have reported tumor-specific
error alongside the domain-average metric.
\vspace{1em}}}
\end{center}
\vspace{0.3em}
\begin{center}
\textbf{Fig. 25.} The reconstructed results with varying 3, 6
and 9 surface displacement datasets corresponding to Sample 1
(5\,\% noise is introduced into the measured data, Unit: MPa).
\end{center}

\vspace{1em}

% ---------- Discussion paragraphs ----------
\noindent
In the numerical cases presented in the paper, the total number
of predefined interfaces (for layered structures) or components
(for composites with embedded inclusions) in the initial guesses
is intentionally kept not much larger than the actual target.
This is because the topological information of the elastic
property distribution can often be determined beforehand based
on prior knowledge or observations. For example, in biomedical
imaging, characteristics such as the number, location, and rough
shape of tumors can be observed and identified. Similarly, for
typical layered biological tissues like skins or corneas, the
total number of layers has likely been well studied through in
vitro experiments or other research. Due to this available prior
knowledge, it is unnecessary to set a large number of
deformable components or interfaces in the initial guess for
the explicit inverse method. By constraining the initial guesses
to a reasonable and relevant number of interfaces or inclusions,
the optimization process can be more focused and efficient.
This can lead to faster convergence and more accurate
characterization of the elastic property distribution with a
reduced risk of convergence to a non-physical solution.

However, it should be noted that the results of this study do
not conclusively demonstrate that the explicit inverse method
outperforms the implicit inverse method in all cases. In
scenarios where the elastic property distribution lacks
significant heterogeneity or where no prior knowledge is
available, the explicit inverse method may indeed perform less
effectively compared to the implicit inverse method. This is
because the explicit inverse method relies on predefined
interfaces or inclusions, which may not accurately capture the
complexity of the underlying distribution. In such situations,
the implicit inverse method, with its greater flexibility and
ability to explore a wider range of possibilities, might yield
more accurate and robust results. The choice between the
explicit and implicit inverse methods, therefore, depends on
the specific characteristics of the problem and the availability
of prior information, ensuring the selection of the most
suitable approach for achieving meaningful and reliable results.

Due to the limited availability of experimental datasets, this
paper primarily focuses on numerical verification. It should be
noted that the implicit inverse method has demonstrated
successful applications in reconstructing the nonhomogeneous
shear modulus distribution of soft solids using experimental
datasets [38]. However, considering the superior performance of
the explicit inverse solver in numerical validation, we are
confident that it can achieve even better results when applied
to the experimental datasets. In

% [page ends here — continues on page 20]

\newpage

% =========== Page 20: wei_page20.tex ===========

% ---------- Figure 26 placeholder ----------
\noindent
\begin{center}
\fbox{\parbox{0.92\textwidth}{\centering\vspace{1em}
\textbf{[Figure 26 --- placeholder]}\\[0.8em]
Six-panel comparison of \textbf{Explicit Inverse Method} (top
row) vs.\ \textbf{Implicit Inverse Method} (bottom row) on
\textbf{Case 3 Sample 2} (two embedded inclusions side-by-side,
target $\approx 5$\,MPa each), with \textbf{no noise}, using 3,
6, and 9 displacement datasets.\\[0.4em]
\textbf{Explicit (top row)} --- both inclusions are recovered as
distinct red blobs on a uniform blue background. Colorbar
ranges:
(a) 3 fields, SM $\in [1.00, 7.06]$\,MPa --- \textbf{41\,\%
overshoot};
(b) 6 fields, SM $\in [1.00, 6.06]$\,MPa --- 21\,\% overshoot;
(c) 9 fields, SM $\in [1.00, 5.97]$\,MPa --- 19\,\% overshoot.
Max drifts downward monotonically as more fields are added,
approaching but not reaching the 5\,MPa nominal.\\[0.4em]
\textbf{Implicit (bottom row)} --- both inclusions are
recovered as green/yellow blobs (distinctly underbright relative
to the explicit method's red). Colorbar ranges:
(d) 3 fields, SM $\in [1.00, 2.61]$\,MPa --- \textbf{48\,\%
undershoot};
(e) 6 fields, SM $\in [1.00, 3.00]$\,MPa --- 40\,\% undershoot;
(f) 9 fields, SM $\in [1.00, 3.82]$\,MPa --- 24\,\% undershoot.
Max drifts upward monotonically, approaching but nowhere close
to the 5\,MPa nominal.\\[0.3em]
\textbf{Takeaway}: on the two-inclusion case (even without
noise), the explicit method overshoots by 19--41\,\% and the
implicit method undershoots by 24--48\,\%. The paper's prose
reports both deviations (see page~18: ``overestimated by
approximately 20\,\%'' for explicit, ``underestimation ...\ by
at least 20\,\%'' for implicit). In this case the prose framing
and the visible data agree --- unlike the Fig.~25 case where
the ``2\,\% average error'' claim hid a 7\,\% inclusion-specific
undershoot.
\vspace{1em}}}
\end{center}
\vspace{0.3em}
\begin{center}
\textbf{Fig. 26.} The reconstructed results with varying 3, 6
and 9 surface displacement datasets corresponding to Sample 2
(Without noise, Unit: MPa).
\end{center}

\vspace{1em}

% ---------- Table 6 ----------
\noindent\textbf{Table 6}\\[0.2em]
{\small The regularization factors of Figs.~23--26.}
\vspace{0.3em}

\begin{center}
\footnotesize
\begin{tabular}{@{}ll ccc ccc@{}}
\toprule
\multicolumn{2}{l}{Inverse method}
  & \multicolumn{3}{c}{Explicit inverse method}
  & \multicolumn{3}{c}{Implicit inverse method} \\
\cmidrule(lr){3-5}\cmidrule(lr){6-8}
\multicolumn{2}{l}{Number of surface displacement datasets}
  & 3 & 6 & 9 & 3 & 6 & 9 \\
\midrule
\multirow{3}{*}{Sample 1}
  & no noise    & $1\times 10^{-15}$ & $1\times 10^{-16}$ & $1\times 10^{-17}$ & $1\times 10^{-14}$ & $1\times 10^{-14}$ & $1\times 10^{-14}$ \\
  & 1\,\% noise & $1\times 10^{-13}$ & $1\times 10^{-13}$ & $2\times 10^{-13}$ & $5\times 10^{-13}$ & $3\times 10^{-13}$ & $2\times 10^{-12}$ \\
  & 5\,\% noise & $1\times 10^{-12}$ & $1\times 10^{-12}$ & $1\times 10^{-12}$ & $1\times 10^{-11}$ & $1\times 10^{-11}$ & $1\times 10^{-12}$ \\
\midrule
Sample 2 & no noise & 0.00 & 0.00 & 0.00 & $1\times 10^{-13}$ & $1\times 10^{-13}$ & $1\times 10^{-14}$ \\
\bottomrule
\end{tabular}
\end{center}

\vspace{1em}

% ---------- Prose transition ----------
\noindent
future research, we plan to explore the integration of optical
measurement techniques with the proposed explicit inverse
approach to conduct practical experiments and further test this
novel structural health monitoring approach. By combining
optical measurement data with our method, we aim to validate
its effectiveness in real-world scenarios and enhance its
applicability for monitoring the variation of biomechanical
behavior of soft tissues. This would provide valuable insights
into the method's performance under realistic conditions and
potentially open up new avenues for its practical implementation
in biomedical applications.

Overall, the explicit inverse approach presented in the paper
offers a more efficient and cost-effective method for
characterizing the elastic property distribution of biological
structures with surface deformation. The reduction in
optimization variables and the reliance on easily measurable
surface deformation data make it a valuable contribution to the
field.

\vspace{0.5em}

\noindent\textbf{Conclusion}

\vspace{0.3em}

\noindent
In this paper, we propose an explicit inverse approach to map
the nonhomogeneous elastic property distribution of three
typical tissue structure by surface deformation nondestructively.
Several numerical examples are employed to assess the
feasibility and performance of the proposed method. The results
of the numerical study demonstrate the robust capability of the
explicit inverse method in accurately reconstructing the
nonhomogeneous shear modulus distribution of soft tissues,
surpassing the performance of the implicit inverse method with
limited surface deformation datasets. Specifically for the
layered ring structure, we observe that the average relative
error of the reconstruction can exceed 40\,\% when using the
implicit inverse method, while the explicit inverse method
achieves a remarkable average relative error of only 5\,\%.
The outcomes of our research underscore the efficacy of the
proposed explicit inverse approach, offering a more efficient
and cost-effective means to characterize the elastic property
distribution of biological structures utilizing a reduced
number of optimization variables.

% [page ends here — continues on page 21]

\newpage

% =========== Page 21: wei_page21.tex ===========

% ---------- Figure 27 placeholder ----------
\noindent
\begin{center}
\fbox{\parbox{0.92\textwidth}{\centering\vspace{1em}
\textbf{[Figure 27 --- placeholder]}\\[0.8em]
Six-panel comparison of \textbf{two different inclusion shapes}
reconstructed with 3 displacement datasets (Sample 1, no noise).
The two rows show different target shapes; the columns show
target, explicit method result, and implicit method result.\\[0.4em]
\textbf{Top row (irregular/elongated inclusion)}:
(a) target: red elongated/irregular shape, SM $\in [1.00, 5.00]$\,MPa;
(b) Explicit, 3 fields: SM $\in [1.00, 5.00]$\,MPa --- the
reconstructed shape faithfully preserves the irregular target
geometry;
(c) Implicit, 3 fields: SM $\in [1.00, 4.98]$\,MPa --- the shape
is \emph{distorted}, rendered as a green/yellow blob without
the sharp edges of the target, and fails to capture the elongated
profile.\\[0.4em]
\textbf{Bottom row (square inclusion)}:
(d) target: red square shape, SM $\in [1.00, 5.00]$\,MPa;
(e) Explicit, 3 fields: SM $\in [1.00, 5.00]$\,MPa --- the
reconstructed shape is a rounded-square, closely matching the
target;
(f) Implicit, 3 fields: SM $\in [1.00, 4.71]$\,MPa --- the shape
is rendered as a green amorphous blob, significantly distorting
the square target.\\[0.3em]
\textbf{Takeaway}: this figure is the qualitative case for the
explicit inverse method's main selling point on non-circular
inclusion geometries. The explicit method's B-spline-based shape
parameterization allows it to faithfully represent elongated or
square shapes; the implicit method's pixel-based optimization
reverts to blob-like reconstructions regardless of the target
shape. This is a shape-fidelity argument, not a shear-modulus
accuracy argument --- and it's where the explicit method's prior
knowledge (topology + rough shape) pays off most visibly.
\vspace{1em}}}
\end{center}
\vspace{0.3em}
\begin{center}
\textbf{Fig. 27.} The target and reconstructed results for
varying shapes of inclusion with 3 displacement datasets
corresponding to Sample 1 (Without noise, Unit: MPa).
\end{center}

\vspace{1em}

% ---------- CRediT authorship contribution statement ----------
\noindent\textbf{CRediT authorship contribution statement}

\vspace{0.3em}

\noindent
\textbf{Yue Mei:} Writing -- review \& editing, Writing --
original draft, Methodology, Formal analysis, Conceptualization.
\textbf{Dongmei Zhao:} Writing -- review \& editing, Writing --
original draft, Software, Formal analysis, Data curation.
\textbf{Changjiang Xiao:} Writing -- original draft,
Investigation, Formal analysis, Data curation.
\textbf{Zhi Sun:} Methodology, Formal analysis, Data curation.
\textbf{Weisheng Zhang:} Writing -- review \& editing, Writing
-- original draft, Supervision, Funding acquisition.
\textbf{Xu Guo:} Writing -- review \& editing, Writing --
original draft, Supervision, Funding acquisition.

\vspace{0.5em}

% ---------- Declaration of competing interest ----------
\noindent\textbf{Declaration of competing interest}

\vspace{0.3em}

\noindent
The authors declare that they have no known competing financial
interests or personal relationships that could have appeared to
influence the work reported in this paper.

\vspace{0.5em}

% ---------- Data availability ----------
\noindent\textbf{Data availability}

\vspace{0.3em}

\noindent
Data will be made available on request.

\vspace{0.5em}

% ---------- Acknowledgments ----------
\noindent\textbf{Acknowledgments}

\vspace{0.3em}

\noindent
The authors acknowledge the support from Innovative Research
Groups of the National Natural Science Foundation of China
(11821202), National Natural Science Foundation of China
(12002075), Liao Ning Revitalization Talents Program
(XLYC1907119), Program for Changjiang Scholars, Innovative
Research Team in University (PCSIRT) and 111 Project (B14013).

\noindent
We also acknowledge Bo Wang and Jiaqi Hong for their helpful
suggestions on the revised manuscript.

\vspace{0.5em}

% ---------- References [1]-[4] ----------
\noindent\textbf{References}

\vspace{0.3em}

\begin{list}{}{%
  \setlength{\leftmargin}{2em}%
  \setlength{\itemindent}{-2em}%
  \setlength{\itemsep}{0.5em}%
  \setlength{\parsep}{0pt}%
}

\item [1] O. Osanai, M. Ohtsuka, M. Hotta, T. Kitaharai, Y.
Takema, A new method for the visualization and quantification
of internal skin elasticity by ultrasound imaging,
\emph{Skin Res. Technol.} 17 (3) (2011) 270--277,
https://doi.org/10.1111/j.1600-0846.2010.00492.x.

\item [2] Y. Zhao, B. Feng, J. Lee, N. Lu, D.M. Pierce, A
multi-layered computational model for wrinkling of human skin
predicts aging effects, \emph{J. Mech. Behav. Biomed. Mater.}
103 (2020) 103552, https://doi.org/10.1016/j.jmbbm.2019.103552.

\item [3] I.L. Kruglikov, P.E. Scherer, Skin aging as a
mechanical phenomenon: the main weak links, \emph{Nutr. Healthy
Aging} 4 (4) (2018) 291--307,
https://doi.org/10.3233/NHA-170037.

\item [4] G.A. Holzapfel, T.C. Gasser, R.W. Ogden, A New
Constitutive Framework for Arterial Wall Mechanics and a
Comparative Study of Material Models, \emph{J. Elast. Phys.
Sci. Solids} 61 (1) (2000) 1--48,
https://doi.org/10.1023/A:1010835316564.

\end{list}

% [page ends here — continues on page 22]

\newpage

% =========== Page 22: wei_page22.tex ===========

\begin{list}{}{%
  \setlength{\leftmargin}{2em}%
  \setlength{\itemindent}{-2em}%
  \setlength{\itemsep}{0.5em}%
  \setlength{\parsep}{0pt}%
}

\item [5] N.J.B. Driessen, W. Wilson, C.V.C. Bouten, F.P.T.
Baaijens, A computational model for collagen fibre remodelling
in the arterial wall, \emph{J. Theor. Biol.} 226 (1) (2004)
53--64, https://doi.org/10.1016/j.jtbi.2003.08.004.

\item [6] F. Zvietcovich, P. Pongchalee, P. Meemon, J.P.
Rolland, K.J. Parker, Reverberant 3D optical coherence
elastography maps the elasticity of individual corneal layers,
\emph{Nat. Commun.} 10 (1) (2019) 4895,
https://doi.org/10.1038/s41467-019-12803-4.

\item [7] J. Ma, Y. Wang, P. Wei, V. Jhanji, Biomechanics and
structure of the cornea: implications and association with
corneal disorders, \emph{Surv. Ophthalmol.} 63 (6) (2018)
851--861, https://doi.org/10.1016/j.survophthal.2018.05.004.

\item [8] S. Goenezen, J. Dord, Z. Sink, P.E. Barbone, J. Jiang,
T.J. Hall, A.A. Oberai, Linear and nonlinear elastic modulus
imaging: an application to breast cancer diagnosis, \emph{IEEE
Trans. Med. Imaging} 31 (8) (2012) 1628--1637,
https://doi.org/10.1109/TMI.2012.2201497.

\item [9] B. Coudrillier, J. Tian, S. Alexander, K.M. Myers,
H.A. Quigley, T.D. Nguyen, Biomechanics of the Human Posterior
Sclera: age- and Glaucoma-Related Changes Measured Using
Inflation Testing, \emph{Invest. Ophthalmol. Vis. Sci.} 53 (4)
(2012) 1714--1728, https://doi.org/10.1167/iovs.11-8009.

\item [10] M. Butlin, A.P. Avolio, Age-Related Changes in the
Mechanical Properties of Large Arteries, in: B. Derby, K.
Akhtar (Eds.), \emph{Mechanical Properties of Aging Soft
Tissues}, Springer International Publishing, Cham, 2015,
pp.~37--74, https://doi.org/10.1007/978-3-319-03970-1\_3.

\item [11] P.G. Agache, C. Monneur, J.L. Leveque, J. De Rigal,
Mechanical properties and Young's modulus of human skin in
vivo, \emph{Arch. Dermatol. Res.} 269 (3) (1980) 221--232,
https://doi.org/10.1007/BF00406415.

\item [12] I.Z. Apostolakis, G.M. Karageorgos, P. Nauleau, S.D.
Nandlall, E.E. Konofagou, Adaptive Pulse Wave Imaging:
automated Spatial Vessel Wall Inhomogeneity Detection in
Phantoms and in-Vivo, \emph{IEEE Trans. Med. Imaging} 39 (1)
(2020) 259--269, https://doi.org/10.1109/TMI.2019.2926141.

\item [13] M.B. Bersi, et al., Multimodality Imaging-Based
Characterization of Regional Material Properties in a Murine
Model of Aortic Dissection, \emph{Sci. Rep.} 10 (1) (2020)
9244, https://doi.org/10.1038/s41598-020-65624-7.

\item [14] M. Ambroziak, P. Pietruski, B. Noszczyk, L. Paluch,
Ultrasonographic elastography in the evaluation of normal and
pathological skin -- a review, \emph{Postepy Dermatol. Alergol.}
36 (6) (2019) 667--672, https://doi.org/10.5114/ada.2018.77069.

\item [15] C. Pailler-Mattei, S. Bec, H. Zahouani, In vivo
measurements of the elastic mechanical properties of human
skin by indentation tests, \emph{Med. Eng. Phys.} 30 (5) (2008)
599--606, https://doi.org/10.1016/j.medengphy.2007.06.011.

\item [16] G.Y. Li, et al., Mechanical characterization of
functionally graded soft materials with ultrasound elastography,
\emph{Philos. Trans. R. Soc. Math. Phys. Eng. Sci.} 377 (2144)
(2019) 20180075, https://doi.org/10.1098/rsta.2018.0075.

\item [17] G.Y. Li, Y. Zheng, Y.X. Jiang, Z. Zhang, Y. Cao,
Guided wave elastography of layered soft tissues, \emph{Acta
Biomater} 84 (2019) 293--304,
https://doi.org/10.1016/j.actbio.2018.12.002.

\item [18] M. Hajhashemkhani, M. Hematiyan, An inverse method
for elastic constants identification of two-layer hyperelastic
bodies with suction loading, \emph{Proc. Inst. Mech. Eng., Part L}
(2023), https://doi.org/10.1177/14644207231203294.

\item [19] M. Hajhashemkhani, M. Hematiyan, S. Goenezen, Inverse
determination of elastic constants of a hyper-elastic member
with inclusions using simple displacement/length measurements,
\emph{J. Strain Anal. Eng. Des.} 53 (7) (2018) 529--542,
https://doi.org/10.1177/0309324718792452.

\item [20] Z. Liu, Y. Sun, J. Deng, D. Zhao, Y. Mei, J. Luo, A
Comparative Study of Direct and Iterative Inversion Approaches
to Determine the Spatial Shear Modulus Distribution of Elastic
Solids, \emph{Int. J. Appl. Mech.} 11 (10) (2019) 1950097,
https://doi.org/10.1142/S1758825119500977.

\item [21] Y. Mei, S. Goenezen, Quantifying the anisotropic
linear elastic behavior of solids, \emph{Int. J. Mech. Sci.}
163 (2019) 105131,
https://doi.org/10.1016/j.ijmecsci.2019.105131.

\item [22] Y. Mei, et al., A comparative study of two
constitutive models within an inverse approach to determine
the spatial stiffness distribution in soft materials,
\emph{Int. J. Mech. Sci.} 140 (2018) 446--454,
https://doi.org/10.1016/j.ijmecsci.2018.03.004.

\item [23] Z. Zhao, et al., Optical Coherence Elastography of
3D bilayer soft solids using full-field and partial displacement
measurements, \emph{Med. Nov. Technol. Devices} (2022) 100134,
https://doi.org/10.1016/j.medndl.2022.100134.

\item [24] D.I. Grodin, et al., Repeatability of Linear and
Nonlinear Elastic Modulus Maps From Repeat Scans in the Breast,
\emph{IEEE Trans. Med. Imaging} 40 (2) (Feb.~2021) 748--757,
https://doi.org/10.1109/TMI.2020.3036552.

\item [25] Y. Mei, S. Wang, X. Shen, S. Rabke, S. Goenezen,
Mechanics Based Topography: a Preliminary Feasibility Study,
\emph{Sensors} 17 (5) (2017),
https://doi.org/10.3390/s17051075.

\item [26] Y. Mei, R. Fulmer, V. Raja, S. Wang, S. Goenezen,
Estimating the non-homogeneous elastic modulus distribution
from surface deformations, \emph{Int. J. Solids Struct.} 83
(2016) 73--80, https://doi.org/10.1016/j.ijsolstr.2016.01.001.

\item [27] Y. Mei, Z. Du, D. Zhao, W. Zhang, C. Liu, X. Guo,
Moving Morphable Inclusion Approach: an Explicit Framework to
Solve Inverse Problem in Elasticity, \emph{J. Appl. Mech.} 88
(4) (2020), https://doi.org/10.1115/1.4049142.

\item [28] D. Liu, J. Du, A Moving Morphable Components Based
Shape Reconstruction Framework for Electrical Capacitance
Tomography, \emph{IEEE Trans. Med. Imaging} 38 (12) (2019)
2937--2948, https://doi.org/10.1109/TMI.2019.2918566.

\item [29] D. Liu, D. Gu, D. Smyl, J. Deng, J. Du, Shape
Reconstruction Using Boolean Operations in Electrical
Impedance Tomography, \emph{IEEE Trans. Med. Imaging} 39 (9)
(2020) 2954--2964, https://doi.org/10.1109/TMI.2020.2983555.

\item [30] T.J.R. Hughes, L.P. Franca, M. Balestra, A new
finite element formulation for computational fluid dynamics:
V. Circumventing the Babuska-Brezzi condition: a stable
Petrov-Galerkin formulation of the Stokes problem
accommodating equal-order interpolations, \emph{Comput. Methods
Appl. Mech. Engrg.} 59 (1) (1986) 85--99,
https://doi.org/10.1016/0045-7825(86)90025-3.

\item [31] W. Zhang, Z. Song, J. Zhou, Z. Du, Y. Zhu, Z. Sun,
X. Guo, Topology optimization with multiple materials via
moving morphable component (MMC) method, \emph{Int. J. Numer.
Methods Eng.} 133 (11) (2018) 1653--1675,
https://doi.org/10.1002/nme.5714.

\item [32] J.N. Reddy, \emph{Introduction to the Finite Element
Method}, McGraw-Hill Education, 2019.

\item [33] Y. Mei, S. Goenezen, Spatially Weighted Objective
Function to Solve the Inverse Elasticity Problem for the
Elastic Modulus, in: B. Doyle, K. Miller, A. Wittek, P. Nielsen
(Eds.), \emph{Computational Biomechanics for Medicine},
Springer International Publishing, Cham, 2015, pp.~47--58,
https://doi.org/10.1007/978-3-319-15503-6\_5.

\item [34] A.A. Oberai, N.H. Gokhale, G.R. Feij\'oo, Solution of
inverse problems in elasticity imaging using the adjoint
method, \emph{Inverse Probl} 19 (2) (2003) 297,
https://doi.org/10.1088/0266-5611/19/2/304.

\item [35] S. Goenezen, P. Barbone, A.A. Oberai, Solution of
the nonlinear elasticity imaging inverse problem: the
incompressible case, \emph{Comput. Method. Appl. M.} 200
(13--16) (2011) 1406--1420,
https://doi.org/10.1016/j.cma.2010.12.018.

\item [36] A. Kalra, A. Lowe, A.M. Al-Jumaily, Mechanical
behaviour of skin: a review, \emph{J. Mater. Sci. Eng.} 5
(2016) 1--7, https://doi.org/10.4172/2169-0022.1000254.

\item [37] F. Gao, Z. Guo, M. Sakamoto, et al., Fluid-structure
Interaction within a Layered Aortic Wall Model, \emph{J. Biol.
Phys.} 32 (2006) 435--454,
https://doi.org/10.1007/s10867-006-9027-7.

\item [38] S. Goenezen, B.J. Kim, M. Kotecha, P. Luo, M.R.
Hematiyan, Mechanics based tomography (MBT): validation using
experimental data, \emph{J. Mech. Phys. Solids} 146 (2021)
104187, https://doi.org/10.1016/j.jmps.2020.104187.

\end{list}

% [end of paper — page 22 of 22]

\end{document}
